%% bare_conf.tex
%% V1.4b
%% 2015/08/26
%% by Michael Shell
%% See:
%% http://www.michaelshell.org/
%% for current contact information.
%%
%% This is a skeleton file demonstrating the use of IEEEtran.cls
%% (requires IEEEtran.cls version 1.8b or later) with an IEEE
%% conference paper.
%%
%% Support sites:
%% http://www.michaelshell.org/tex/ieeetran/
%% http://www.ctan.org/pkg/ieeetran
%% and
%% http://www.ieee.org/

%%*************************************************************************
%% Legal Notice:
%% This code is offered as-is without any warranty either expressed or
%% implied; without even the implied warranty of MERCHANTABILITY or
%% FITNESS FOR A PARTICULAR PURPOSE! 
%% User assumes all risk.
%% In no event shall the IEEE or any contributor to this code be liable for
%% any damages or losses, including, but not limited to, incidental,
%% consequential, or any other damages, resulting from the use or misuse
%% of any information contained here.
%%
%% All comments are the opinions of their respective authors and are not
%% necessarily endorsed by the IEEE.
%%
%% This work is distributed under the LaTeX Project Public License (LPPL)
%% ( http://www.latex-project.org/ ) version 1.3, and may be freely used,
%% distributed and modified. A copy of the LPPL, version 1.3, is included
%% in the base LaTeX documentation of all distributions of LaTeX released
%% 2003/12/01 or later.
%% Retain all contribution notices and credits.
%% ** Modified files should be clearly indicated as such, including  **
%% ** renaming them and changing author support contact information. **
%%*************************************************************************


% *** Authors should verify (and, if needed, correct) their LaTeX system  ***
% *** with the testflow diagnostic prior to trusting their LaTeX platform ***
% *** with production work. The IEEE's font choices and paper sizes can   ***
% *** trigger bugs that do not appear when using other class files.       ***                          ***
% The testflow support page is at:
% http://www.michaelshell.org/tex/testflow/



\documentclass[conference]{IEEEtran}
% Some Computer Society conferences also require the compsoc mode option,
% but others use the standard conference format.
%
% If IEEEtran.cls has not been installed into the LaTeX system files,
% manually specify the path to it like:
% \documentclass[conference]{../sty/IEEEtran}





% Some very useful LaTeX packages include:
% (uncomment the ones you want to load)


% *** MISC UTILITY PACKAGES ***
%
%\usepackage{ifpdf}
% Heiko Oberdiek's ifpdf.sty is very useful if you need conditional
% compilation based on whether the output is pdf or dvi.
% usage:
% \ifpdf
%   % pdf code
% \else
%   % dvi code
% \fi
% The latest version of ifpdf.sty can be obtained from:
% http://www.ctan.org/pkg/ifpdf
% Also, note that IEEEtran.cls V1.7 and later provides a builtin
% \ifCLASSINFOpdf conditional that works the same way.
% When switching from latex to pdflatex and vice-versa, the compiler may
% have to be run twice to clear warning/error messages.






% *** CITATION PACKAGES ***
%
%\usepackage{cite}
% cite.sty was written by Donald Arseneau
% V1.6 and later of IEEEtran pre-defines the format of the cite.sty package
% \cite{} output to follow that of the IEEE. Loading the cite package will
% result in citation numbers being automatically sorted and properly
% "compressed/ranged". e.g., [1], [9], [2], [7], [5], [6] without using
% cite.sty will become [1], [2], [5]--[7], [9] using cite.sty. cite.sty's
% \cite will automatically add leading space, if needed. Use cite.sty's
% noadjust option (cite.sty V3.8 and later) if you want to turn this off
% such as if a citation ever needs to be enclosed in parenthesis.
% cite.sty is already installed on most LaTeX systems. Be sure and use
% version 5.0 (2009-03-20) and later if using hyperref.sty.
% The latest version can be obtained at:
% http://www.ctan.org/pkg/cite
% The documentation is contained in the cite.sty file itself.






% *** GRAPHICS RELATED PACKAGES ***
%
\ifCLASSINFOpdf
  % \usepackage[pdftex]{graphicx}
  % declare the path(s) where your graphic files are
  % \graphicspath{{../pdf/}{../jpeg/}}
  % and their extensions so you won't have to specify these with
  % every instance of \includegraphics
  % \DeclareGraphicsExtensions{.pdf,.jpeg,.png}
\else
  % or other class option (dvipsone, dvipdf, if not using dvips). graphicx
  % will default to the driver specified in the system graphics.cfg if no
  % driver is specified.
  % \usepackage[dvips]{graphicx}
  % declare the path(s) where your graphic files are
  % \graphicspath{{../eps/}}
  % and their extensions so you won't have to specify these with
  % every instance of \includegraphics
  % \DeclareGraphicsExtensions{.eps}
\fi
% graphicx was written by David Carlisle and Sebastian Rahtz. It is
% required if you want graphics, photos, etc. graphicx.sty is already
% installed on most LaTeX systems. The latest version and documentation
% can be obtained at: 
% http://www.ctan.org/pkg/graphicx
% Another good source of documentation is "Using Imported Graphics in
% LaTeX2e" by Keith Reckdahl which can be found at:
% http://www.ctan.org/pkg/epslatex
%
% latex, and pdflatex in dvi mode, support graphics in encapsulated
% postscript (.eps) format. pdflatex in pdf mode supports graphics
% in .pdf, .jpeg, .png and .mps (metapost) formats. Users should ensure
% that all non-photo figures use a vector format (.eps, .pdf, .mps) and
% not a bitmapped formats (.jpeg, .png). The IEEE frowns on bitmapped formats
% which can result in "jaggedy"/blurry rendering of lines and letters as
% well as large increases in file sizes.
%
% You can find documentation about the pdfTeX application at:
% http://www.tug.org/applications/pdftex





% *** MATH PACKAGES ***
%
%\usepackage{amsmath}
% A popular package from the American Mathematical Society that provides
% many useful and powerful commands for dealing with mathematics.
%
% Note that the amsmath package sets \interdisplaylinepenalty to 10000
% thus preventing page breaks from occurring within multiline equations. Use:
%\interdisplaylinepenalty=2500
% after loading amsmath to restore such page breaks as IEEEtran.cls normally
% does. amsmath.sty is already installed on most LaTeX systems. The latest
% version and documentation can be obtained at:
% http://www.ctan.org/pkg/amsmath





% *** SPECIALIZED LIST PACKAGES ***
%
%\usepackage{algorithmic}
% algorithmic.sty was written by Peter Williams and Rogerio Brito.
% This package provides an algorithmic environment fo describing algorithms.
% You can use the algorithmic environment in-text or within a figure
% environment to provide for a floating algorithm. Do NOT use the algorithm
% floating environment provided by algorithm.sty (by the same authors) or
% algorithm2e.sty (by Christophe Fiorio) as the IEEE does not use dedicated
% algorithm float types and packages that provide these will not provide
% correct IEEE style captions. The latest version and documentation of
% algorithmic.sty can be obtained at:
% http://www.ctan.org/pkg/algorithms
% Also of interest may be the (relatively newer and more customizable)
% algorithmicx.sty package by Szasz Janos:
% http://www.ctan.org/pkg/algorithmicx




% *** ALIGNMENT PACKAGES ***
%
%\usepackage{array}
% Frank Mittelbach's and David Carlisle's array.sty patches and improves
% the standard LaTeX2e array and tabular environments to provide better
% appearance and additional user controls. As the default LaTeX2e table
% generation code is lacking to the point of almost being broken with
% respect to the quality of the end results, all users are strongly
% advised to use an enhanced (at the very least that provided by array.sty)
% set of table tools. array.sty is already installed on most systems. The
% latest version and documentation can be obtained at:
% http://www.ctan.org/pkg/array


% IEEEtran contains the IEEEeqnarray family of commands that can be used to
% generate multiline equations as well as matrices, tables, etc., of high
% quality.




% *** SUBFIGURE PACKAGES ***
%\ifCLASSOPTIONcompsoc
%  \usepackage[caption=false,font=normalsize,labelfont=sf,textfont=sf]{subfig}
%\else
%  \usepackage[caption=false,font=footnotesize]{subfig}
%\fi
% subfig.sty, written by Steven Douglas Cochran, is the modern replacement
% for subfigure.sty, the latter of which is no longer maintained and is
% incompatible with some LaTeX packages including fixltx2e. However,
% subfig.sty requires and automatically loads Axel Sommerfeldt's caption.sty
% which will override IEEEtran.cls' handling of captions and this will result
% in non-IEEE style figure/table captions. To prevent this problem, be sure
% and invoke subfig.sty's "caption=false" package option (available since
% subfig.sty version 1.3, 2005/06/28) as this is will preserve IEEEtran.cls
% handling of captions.
% Note that the Computer Society format requires a larger sans serif font
% than the serif footnote size font used in traditional IEEE formatting
% and thus the need to invoke different subfig.sty package options depending
% on whether compsoc mode has been enabled.
%
% The latest version and documentation of subfig.sty can be obtained at:
% http://www.ctan.org/pkg/subfig




% *** FLOAT PACKAGES ***
%
%\usepackage{fixltx2e}
% fixltx2e, the successor to the earlier fix2col.sty, was written by
% Frank Mittelbach and David Carlisle. This package corrects a few problems
% in the LaTeX2e kernel, the most notable of which is that in current
% LaTeX2e releases, the ordering of single and double column floats is not
% guaranteed to be preserved. Thus, an unpatched LaTeX2e can allow a
% single column figure to be placed prior to an earlier double column
% figure.
% Be aware that LaTeX2e kernels dated 2015 and later have fixltx2e.sty's
% corrections already built into the system in which case a warning will
% be issued if an attempt is made to load fixltx2e.sty as it is no longer
% needed.
% The latest version and documentation can be found at:
% http://www.ctan.org/pkg/fixltx2e


%\usepackage{stfloats}
% stfloats.sty was written by Sigitas Tolusis. This package gives LaTeX2e
% the ability to do double column floats at the bottom of the page as well
% as the top. (e.g., "\begin{figure*}[!b]" is not normally possible in
% LaTeX2e). It also provides a command:
%\fnbelowfloat
% to enable the placement of footnotes below bottom floats (the standard
% LaTeX2e kernel puts them above bottom floats). This is an invasive package
% which rewrites many portions of the LaTeX2e float routines. It may not work
% with other packages that modify the LaTeX2e float routines. The latest
% version and documentation can be obtained at:
% http://www.ctan.org/pkg/stfloats
% Do not use the stfloats baselinefloat ability as the IEEE does not allow
% \baselineskip to stretch. Authors submitting work to the IEEE should note
% that the IEEE rarely uses double column equations and that authors should try
% to avoid such use. Do not be tempted to use the cuted.sty or midfloat.sty
% packages (also by Sigitas Tolusis) as the IEEE does not format its papers in
% such ways.
% Do not attempt to use stfloats with fixltx2e as they are incompatible.
% Instead, use Morten Hogholm'a dblfloatfix which combines the features
% of both fixltx2e and stfloats:
%
% \usepackage{dblfloatfix}
% The latest version can be found at:
% http://www.ctan.org/pkg/dblfloatfix




% *** PDF, URL AND HYPERLINK PACKAGES ***
%
%\usepackage{url}
% url.sty was written by Donald Arseneau. It provides better support for
% handling and breaking URLs. url.sty is already installed on most LaTeX
% systems. The latest version and documentation can be obtained at:
% http://www.ctan.org/pkg/url
% Basically, \url{my_url_here}.




% *** Do not adjust lengths that control margins, column widths, etc. ***
% *** Do not use packages that alter fonts (such as pslatex).         ***
% There should be no need to do such things with IEEEtran.cls V1.6 and later.
% (Unless specifically asked to do so by the journal or conference you plan
% to submit to, of course. )

\usepackage{microtype}
\usepackage{graphicx}
\usepackage{subfigure}
\usepackage{booktabs} % for professional tables
\usepackage{amsmath} 
\usepackage{multirow}
\usepackage{array}
\usepackage{pifont}
\newcommand{\cmark}{\ding{51}} % ✓
\newcommand{\xmark}{\ding{55}} % ✗

% correct bad hyphenation here
\hyphenation{op-tical net-works semi-conduc-tor}


\begin{document}
%
% paper title
% Titles are generally capitalized except for words such as a, an, and, as,
% at, but, by, for, in, nor, of, on, or, the, to and up, which are usually
% not capitalized unless they are the first or last word of the title.
% Linebreaks \\ can be used within to get better formatting as desired.
% Do not put math or special symbols in the title.
\title{
    FlexLLM: Composable HLS Library for Flexible Hybrid LLM Accelerator Design
    \vspace{-0.3cm}
}


% author names and affiliations
% use a multiple column layout for up to three different
% affiliations
% \author{\IEEEauthorblockN{Jiahao Zhang}
% \IEEEauthorblockA{School of Electrical and\\Computer Engineering\\
% Georgia Institute of Technology\\
% Atlanta, Georgia 30332--0250\\
% Email: http://www.michaelshell.org/contact.html}
% \and
% \IEEEauthorblockN{Zifan }
% \IEEEauthorblockA{Twentieth Century Fox\\
% Springfield, USA\\
% Email: homer@thesimpsons.com}
% \and
% \IEEEauthorblockN{James Kirk\\ and Montgomery Scott}
% \IEEEauthorblockA{Starfleet Academy\\
% San Francisco, California 96678--2391\\
% Telephone: (800) 555--1212\\
% Fax: (888) 555--1212}}

% conference papers do not typically use \thanks and this command
% is locked out in conference mode. If really needed, such as for
% the acknowledgment of grants, issue a \IEEEoverridecommandlockouts
% after \documentclass

% for over three affiliations, or if they all won't fit within the width
% of the page, use this alternative format:
% 
\author{
    \IEEEauthorblockN{
        Jiahao Zhang\IEEEauthorrefmark{1},
        Zifan He\IEEEauthorrefmark{1},
        Nicholas Fraser\IEEEauthorrefmark{2}, 
        Michaela Blott\IEEEauthorrefmark{2},
        Yizhou Sun\IEEEauthorrefmark{1}, 
        Jason Cong\IEEEauthorrefmark{1}
    }
    \IEEEauthorblockA{
        \IEEEauthorrefmark{1}Computer Science, University of California, Los Angeles, California\\ 
    }
    \IEEEauthorblockA{
        \IEEEauthorrefmark{2}AMD, Dublin, Ireland
    }
    \vspace{-1cm}
}





% use for special paper notices
%\IEEEspecialpapernotice{(Invited Paper)}




% make the title area
\maketitle


% As a general rule, do not put math, special symbols or citations
% in the abstract
\begin{abstract}
% We present \textbf{FlexLLM}, a composable High-Level Synthesis (HLS) library for rapid development of LLM accelerators. FlexLLM enables hybrid architectures combining temporal reuse and spatial dataflow, with parameterized modules customized for the prefill and decode stages and a comprehensive quantization suite for hardware-efficient yet accurate deployment. Using FlexLLM, we build a complete inference system for the Llama-3.2 1B model in under two months with only 1K lines of code. The system includes: (1) a stage-customized accelerator with high-accuracy quantization (12.68 Wiki-2 PPL) surpassing SpinQuant baseline, and (2) a Hierarchical Memory Transformer (HMT) plug-in for efficient long-context processing. On the AMD U280 FPGA at 16nm, the accelerator achieves 1.29$\times$ end-to-end speedup, 1.64$\times$ higher decode throughput, and 3.14$\times$ better energy efficiency than an NVIDIA A100 GPU (7nm) running BF16 inference; projected results on the V80 FPGA at 7nm reach 4.71$\times$, 6.55$\times$, and 4.13$\times$, respectively. In long-context scenarios, integrating the HMT plug-in reduces prefill latency by 23.23$\times$ and extends the context window by 64$\times$, delivering 1.10$\times$/4.86$\times$ lower end-to-end latency and 5.21$\times$/6.27$\times$ higher energy efficiency on the U280/V80 compared to the A100 baseline. FlexLLM thus bridges algorithmic innovation in LLM inference and high-performance accelerators with minimal manual effort.

We present FlexLLM, a composable High-Level Synthesis (HLS) library for rapid development of domain-specific LLM accelerators. FlexLLM exposes key architectural degrees of freedom for stage-customized inference, enabling hybrid designs that tailor temporal reuse and spatial dataflow differently for prefill and decode, and provides a comprehensive quantization suite to support accurate low-bit deployment. Using FlexLLM, we build a complete inference system for the Llama-3.2 1B model in under two months with only 1K lines of code. The system includes: (1) a stage-customized accelerator with hardware-efficient quantization (12.68 WikiText-2 PPL) surpassing SpinQuant baseline, and (2) a Hierarchical Memory Transformer (HMT) plug-in for efficient long-context processing. On the AMD U280 FPGA at 16nm, the accelerator achieves 1.29$\times$ end-to-end speedup, 1.64$\times$ higher decode throughput, and 3.14$\times$ better energy efficiency than an NVIDIA A100 GPU (7nm) running BF16 inference; projected results on the V80 FPGA at 7nm reach 4.71$\times$, 6.55$\times$, and 4.13$\times$, respectively. In long-context scenarios, integrating the HMT plug-in reduces prefill latency by 23.23$\times$ and extends the context window by 64$\times$, delivering 1.10$\times$/4.86$\times$ lower end-to-end latency and 5.21$\times$/6.27$\times$ higher energy efficiency on the U280/V80 compared to the A100 baseline. FlexLLM thus bridges algorithmic innovation in LLM inference and high-performance accelerators with minimal manual effort.

\end{abstract}

% no keywords




% For peer review papers, you can put extra information on the cover
% page as needed:
% \ifCLASSOPTIONpeerreview
% \begin{center} \bfseries EDICS Category: 3-BBND \end{center}
% \fi
%
% For peerreview papers, this IEEEtran command inserts a page break and
% creates the second title. It will be ignored for other modes.
\IEEEpeerreviewmaketitle

% \section{Introduction}

% Large Language Models (LLMs) have emerged as one of the most transformative technologies of the modern era, driving breakthroughs across domains such as natural language processing~\cite{devlin2019bert, raffel2020t5}, code synthesis~\cite{chen2021codex, alphacode2022science}, and multimodal understanding~\cite{openai2023gpt4, alayrac2022flamingo}. While cloud providers continue investing heavily in GPU-based infrastructures to support large-scale inference~\cite{kwon2023vllm}, there is a growing demand for \textbf{domain-specific accelerators} that offer customized, efficient, and scalable alternatives for diverse deployment environments—from data centers~\cite{hong2022dfx, chen2024understanding} to edge platforms~\cite{chen2019eyerissv2, huang2025edgellm}.

% Compared to general-purpose CPUs and GPUs, domain-specific accelerators—such as those for CNNs~\cite{zhang2016caffeine, basalama2023flexcnn} or ResNets~\cite{li2017fpga, ma2017end, guan2017fp}—have demonstrated significant improvements in performance and energy efficiency. However, modern LLMs differ substantially from conventional DNNs: they contain billions of parameters and exhibit more complex computation and memory access patterns~\cite{brown2020gpt3, dubey2024llama, liu2024deepseek}, resulting in far higher demands on compute throughput, memory bandwidth, and architectural flexibility. These characteristics present three key challenges in designing efficient LLM accelerators:

% \textbf{Challenge 1: Balancing model accuracy and computation/memory overheads.}  
% Contemporary LLMs contain enormous parameter counts, making model compression indispensable for practical accelerator deployment~\cite{chen2024understanding, zeng2024flightllm}. Quantization offers an effective way to reduce computation and memory cost, but naive or overly aggressive quantization often leads to severe accuracy degradation, limiting usability. For example, applying INT4 SmoothQuant~\cite{xiao2023smoothquant} or GPTQ~\cite{frantar2022gptq} to the Llama~3.2-1B model increases perplexity to over $1e2$ on WikiText-2~\cite{liu2024spinquant}, rendering it impractical. Therefore, supporting advanced and adaptive quantization schemes that balance accuracy and efficiency is critical for domain-specific accelerator design.

% \textbf{Challenge 2: Divergent compute and memory behaviors in different inference stages.}  
% LLM inference, based on the Transformer decoder architecture, typically consists of two distinct stages. During the \emph{prefill stage}, the model processes an entire input prompt in parallel, achieving high arithmetic intensity and abundant parallelism. In contrast, the \emph{decode stage} generates tokens autoregressively, where computation becomes dominated by data dependencies and frequent memory accesses. As a result, the prefill stage is compute-bound, while the decode stage is memory-bandwidth-bound, leading to divergent optimization goals for the same model.

% \begin{figure*}[t]
%     \centering
%     \includegraphics[width=1.0\linewidth]{Figures/Arch_Comparison_new.pdf}
%     \vspace{-0.75cm}
%     \caption{
%         % Comparison of three architectural styles for LLM accelerators across prefill and decode stages. Different colors indicate different tokens, and blocks denote hardware modules. Only the linear layer is shown for simplicity, where $A$ denotes the self-attention module (using Grouped-Query Attention here) and $O$ denotes the output projection layer. (a) Diagram of one Transformer block layer.
%         % (b)(c) \textbf{Temporal architectures} achieve high utilization through module reuse but suffer from frequent off-chip memory access and limited flexibility. 
%         % (d)(e) \textbf{Spatial architectures} map each kernel to dedicated hardware modules, enabling fully on-chip streaming but are sensitive to pipeline stalls caused by data dependencies. 
%         % (f)(g) \textbf{Hybrid architectures} combine temporal reuse with spatial parallelism, enabling flexible pipeline scheduling and stage-specific customization to balance utilization, latency, and flexibility. 
%         Comparison of three architectural styles for LLM accelerators across prefill and decode stages. Different colors indicate different tokens and blocks denote hardware modules. Only the linear layer is shown for clarity, where $A$ denotes the Multi-Head Attention (Grouped-Query Attention here) and $O$ denotes the output projection. 
%         (a) One Transformer block. 
%         (b–c) Temporal architectures achieve high utilization via module reuse but suffer from frequent off-chip memory access and limited flexibility. 
%         (d–e) Spatial architectures dedicate modules per kernel for full on-chip streaming but are sensitive to pipeline stalls. 
%         (f–g) Hybrid architectures combine temporal reuse and spatial parallelism to balance utilization, latency, and flexibility.
%         % \YS{shall we first introduce the backbone model's architecture? The GQA was not obvious to me in the first place, and thus the size of K,V,and Q may confuse other readers. What does O mean? ---- I explain GQA in caption, and plan to simplify this figure, and add a LLM decode layer diagram on the side}
%     }
%     \label{fig:arch_comparison}
%     \vspace{-0.3cm}
% \end{figure*}

% Existing accelerators adopt different architectural paradigms to handle these contrasting behaviors but remain constrained by rigid structures and poor adaptability across stages. \textbf{Temporal architectures}~\cite{hong2022dfx, zeng2024flightllm}, shown as Figure~\ref{fig:arch_comparison}(b)(c), sequentially execute layers on shared compute engines, maximizing utilization but causing frequent off-chip memory accesses in prefill due to insufficient buffering. They also struggle to support heterogeneous precisions or layer types within a unified engine. \textbf{Spatial architectures}~\cite{chen2024allo, ye2025streamtensor}, on the other hand, map each kernel to a dedicated hardware module and stream data through on-chip FIFOs. While this approach achieves high throughput under balanced workloads, it faces resource pressure as model size scales and suffers from pipeline stalls caused by unaligned kernel latencies or intrinsic data dependencies (e.g., LayerNorm reductions or autoregressive attention), as illustrated in Figure~\ref{fig:arch_comparison}(d)(e).  

% \textbf{Challenge 3: High manual design effort for model-specific accelerators.}  
% As new LLM architectures continue to emerge, manually developing model-specific accelerators demands extensive engineering effort. RTL-based design often exceeds 100K lines of code for a single implementation, posing a steep learning curve for ML practitioners. High-Level Synthesis (HLS) frameworks such as Allo~\cite{chen2024allo} and StreamTensor~\cite{ye2025streamtensor} have demonstrated promise by automating RTL generation and MLIR-based mapping, but they largely focus on spatial architectures and offer limited flexibility for constructing hybrid designs that jointly optimize computation and memory efficiency across prefill and decode stages.

% Given these challenges, designing customized accelerators for LLMs is a time-consuming process that often takes 1-2 years, creating a growing mismatch with the rapid pace of LLM innovation, where new generations emerge roughly every six months. To bridge this gap, we develop \textbf{FlexLLM}, a composable HLS library for domain-specific LLM accelerator design. FlexLLM is based on \emph{TAPA}~\cite{guo2023tapa}, enabling \textbf{flexible hybrid architecture construction} that combines the benefits of temporal reuse and spatial parallelism, while integrating advanced quantization techniques for accuracy preservation. The library provides highly parameterized and templated modules optimized for the prefill and decode stages, offering a user-friendly interface for rapid and efficient accelerator customization. We’ll open-source it once the paper is accepted for publication.

% Using FlexLLM, we implement a series of hybrid accelerators for recent LLM works with fewer than 1K lines of C++ code. The implementation seamlessly integrates the Llama-3.2 1B model with state-of-the-art techniques such as SpinQuant~\cite{liu2024spinquant} for low-bit quantization and the Hierarchical Memory Transformer (HMT)~\cite{he2024hmt} plug-in for efficient long-context processing. With FlexLLM, new algorithmic innovations can be quickly incorporated into domain-specific accelerator pipelines, significantly reducing design effort while maintaining high performance and architectural flexibility.

% \noindent In summary, our main contributions are as follows:


% \begin{itemize}[wide=0pt]
%     \item[\textbf{(a)}] We present FlexLLM, the first composable HLS library that supports flexible hybrid architecture construction for LLM accelerators. Its parameterized and templated modules allow users to design model-specific accelerators, significantly reducing design complexity and code volume.
%     \item[\textbf{(b)}] We integrate a comprehensive set of quantization and dequantization modules, enabling accurate and efficient LLM deployment on custom accelerators. To the best of our knowledge, this represents the most advanced quantization support among existing LLM accelerator frameworks.
%     \item[\textbf{(c)}] Leveraging FlexLLM, we construct a stage-customized hybrid accelerator for the Llama-3.2 1B model using a W4A4KV8 SpinQuant scheme. Compared to an NVIDIA A100 80G GPU BF16 baseline, our FPGA implementation achieves 1.29$\times$ end-to-end speedup, 1.64$\times$ higher decode throughput, and 3.14$\times$ better energy efficiency on the AMD Alveo U280, and an estimated 4.71$\times$, 6.55$\times$, and 4.13$\times$ improvements, respectively, on the Versal V80.
%     \item[\textbf{(d)}] We further integrate a HMT plug-in using FlexLLM, reducing prefill latency by up to 23.23$\times$ and extending the model’s effective context window by more than 64$\times$ with less than 7.5\% resource and 0.6\% latency overhead, demonstrating FlexLLM’s scalability for long-context acceleration.
% \end{itemize}

\section{Introduction}

\begin{figure*}[t]
    \centering
    \includegraphics[width=1.0\linewidth]{Figures/Arch_Comparison_new.pdf}
    \vspace{-0.75cm}
    \caption{
        % Comparison of three architectural styles for LLM accelerators across prefill and decode stages. Different colors indicate different tokens, and blocks denote hardware modules. Only the linear layer is shown for simplicity, where $A$ denotes the self-attention module (using Grouped-Query Attention here) and $O$ denotes the output projection layer. (a) Diagram of one Transformer block layer.
        % (b)(c) \textbf{Temporal architectures} achieve high utilization through module reuse but suffer from frequent off-chip memory access and limited flexibility. 
        % (d)(e) \textbf{Spatial architectures} map each kernel to dedicated hardware modules, enabling fully on-chip streaming but are sensitive to pipeline stalls caused by data dependencies. 
        % (f)(g) \textbf{Hybrid architectures} combine temporal reuse with spatial parallelism, enabling flexible pipeline scheduling and stage-specific customization to balance utilization, latency, and flexibility. 
        Comparison of three architectural styles for LLM accelerators across prefill and decode stages. Different colors indicate different tokens and blocks denote hardware modules. Only the linear layer is shown for clarity, where $A$ denotes the Multi-Head Attention (Grouped-Query Attention here) and $O$ denotes the output projection. 
        (a) One Transformer block. 
        (b–c) Temporal architectures achieve high utilization via module reuse but suffer from frequent off-chip memory access and limited flexibility. 
        (d–e) Spatial architectures dedicate modules per kernel for full on-chip streaming but are sensitive to pipeline stalls. 
        (f–g) Hybrid architectures combine temporal reuse and spatial parallelism to balance utilization, latency, and flexibility.
        % \YS{shall we first introduce the backbone model's architecture? The GQA was not obvious to me in the first place, and thus the size of K,V,and Q may confuse other readers. What does O mean? ---- I explain GQA in caption, and plan to simplify this figure, and add a LLM decode layer diagram on the side}
    }
    \label{fig:arch_comparison}
    \vspace{-0.3cm}
\end{figure*}

Large Language Models (LLMs) have emerged as one of the most transformative technologies of the modern era, driving breakthroughs across natural language processing~\cite{devlin2019bert, raffel2020t5}, code generation~\cite{chen2021codex, alphacode2022science}, and multimodal understanding~\cite{openai2023gpt4, alayrac2022flamingo}. While cloud providers continue investing heavily in GPU-based infrastructures to support large-scale inference~\cite{kwon2023vllm}, there is a growing demand for \textbf{domain-specific accelerators} that provide customized, efficient, and scalable alternatives across deployment environments—from data centers~\cite{hong2022dfx, chen2024understanding} to edge platforms~\cite{chen2019eyerissv2, huang2025edgellm}.

Compared to general-purpose CPUs and GPUs, domain-specific accelerators have demonstrated significant gains in performance and energy efficiency for prior DNN workloads (e.g., CNNs~\cite{zhang2016caffeine, basalama2023flexcnn}, ResNets~\cite{li2017fpga, ma2017end, guan2017fp}). However, modern LLMs differ substantially from conventional DNNs: they contain billions of parameters and exhibit more complex computation and memory access patterns~\cite{brown2020gpt3, dubey2024llama, liu2024deepseek}, resulting in far higher demands on compute throughput, memory bandwidth, and architectural flexibility. These characteristics present three key challenges in designing efficient LLM accelerators:

\textbf{Challenge 1: Divergent compute and memory behaviors in prefill vs.\ decode.}  
LLM inference with Transformer decoders consists of two stages with fundamentally different bottlenecks. During \emph{prefill}, the model processes the prompt in parallel, offering high arithmetic intensity and abundant parallelism. During \emph{decode}, tokens are generated autoregressively; computation is dominated by data dependencies and frequent memory accesses, making it strongly memory-bandwidth-bound. As a result, the same model is often compute-bound in prefill but memory-/dependency-bound in decode, creating conflicting optimization goals within one serving pipeline.

Existing FPGA accelerators primarily follow either \textbf{temporal architectures}~\cite{hong2022dfx, zeng2024flightllm} or \textbf{spatial architectures}~\cite{blott2018finn, fastml_hls4ml, chen2024allo, ye2025streamtensor}. Temporal designs reuse shared compute engines across layers, but incur frequent off-chip traffic in prefill due to limited buffering and struggle to support heterogeneous kernels/precisions within a single engine (Fig.~\ref{fig:arch_comparison}(b)(c)). Spatial designs map kernels to dedicated modules and stream intermediate data through on-chip FIFOs, but suffer from pipeline stalls when kernel latencies are unbalanced or when intrinsic dependencies dominate (Fig.~\ref{fig:arch_comparison}(d)(e)). Despite their differences, both paradigms typically share the same implicit choice: using a \emph{single unified architecture} to serve both stages.

However, since prefill and decode are bottlenecked by different resources, a unified (temporal, spatial, or even hybrid) design inevitably over-optimizes one stage while under-serving the other. In other words, \textbf{\textit{any ``one-size-fits-all'' architecture is fundamentally mismatched to LLM serving: prefill and decode must be stage-customized}}, with each stage adopting a different mix of spatial parallelism and temporal reuse to match its distinct compute/memory constraints.


\textbf{Challenge 2: Balancing model accuracy and low-bit compute/memory efficiency.}  
LLMs have enormous parameter counts, making compression indispensable for accelerator deployment~\cite{chen2024understanding, zeng2024flightllm}. Quantization is an effective way to reduce compute and memory cost; however, prior FPGA LLM accelerators often rely on naive integer quantization, which can severely degrade accuracy when pushed to the aggressive low-bit regime required to approach GPU-level throughput, making such designs impractical for real deployment. For example, applying INT4 SmoothQuant~\cite{xiao2023smoothquant} or GPTQ~\cite{frantar2022gptq} to Llama~3.2-1B increases perplexity to over $1e2$ on WikiText-2~\cite{liu2024spinquant}. Given FPGAs' limited compute capability relative to GPUs, \textbf{\textit{practical LLM acceleration must rely on state-of-the-art quantization with hardware-oriented co-design}}, rather than naive quantization pipelines.


\textbf{Challenge 3: High manual effort and slow iteration for model-specific accelerators.}  
As LLM architectures evolve rapidly, developing model-specific accelerators remains highly labor-intensive. RTL-based implementations can exceed 100K lines of code and impose a steep learning curve, making it difficult to iterate at the pace of LLM innovation. Recent HLS frameworks such as Allo~\cite{chen2024allo} and StreamTensor~\cite{ye2025streamtensor} reduce RTL burden via MLIR-based mapping and automated compilation, but they largely target spatial architectures and offer limited support for hybrid designs that optimize prefill and decode specifically.

To close this productivity gap, \textbf{\textit{a high-level programming framework that enables flexible architecture design is necessary for the rapid growth of LLM-oriented accelerators.}} Without a composable framework, accelerator development cycles (often 1--2 years) cannot keep pace with the rapid evolution of LLMs, where new model releases occur on the order of months. What is needed is a framework that exposes the key architectural degrees of freedom—especially stage-customized hybrid design and advanced quantization—while remaining accessible to both ML and accelerator developers.


Motivated by these challenges, we develop \textbf{FlexLLM}, a composable HLS library for domain-specific LLM accelerator design. Built on TAPA~\cite{guo2023tapa}, FlexLLM provides highly parameterized and templated modules that enable rapid construction of hybrid accelerators. Importantly, FlexLLM is, to our knowledge, \textbf{\textit{the first to explicitly go beyond the unified-design paradigm and enable stage-customized hybrid architectures}}. This significantly improves design flexibility and accelerator performance while offering a reusable methodology for future LLM accelerator development. Using FlexLLM, we build a high-performance accelerator system that integrates state-of-the-art LLM techniques with fewer than 1K lines of C++ code, demonstrating that FlexLLM can rapidly incorporate new algorithmic innovations while significantly reducing development effort. We will open-source FlexLLM upon publication.   

In summary, our main contributions include:
\begin{itemize}
    \item We present FlexLLM, a composable HLS library that supports \textit{flexible hybrid, stage-customized architecture construction} for LLM accelerators. Its highly templated modules enable rapid model-specific customization while significantly reducing design complexity and code volume.
    % \item We integrate a comprehensive set of quantization modules, including dynamic/static variants with different symmetric/granularity, and outlier handing modules, enabling accurate and efficient LLM deployment on custom acceleration platforms. To the best of our knowledge, this represents the most advanced quantization support among existing LLM accelerator frameworks.
    \item We integrate a comprehensive quantization stack---covering dynamic and static variants, multiple symmetry and granularity options, and outlier-handling modules---enabling accurate and efficient LLM deployment on customized accelerator platforms. To the best of our knowledge, this provides the most advanced quantization support among existing LLM accelerator frameworks.
    % \item Leveraging FlexLLM, we construct a stage-customized hybrid accelerator for Llama-3.2~1B with a hardware-aware, co-designed W4A4KV8 SpinQuant~\cite{liu2024spinquant} scheme, achieving 12.68 WikiText-2 PPL better than the SpinQuant baseline. Compared to the BF16 baseline on an NVIDIA A100 GPU, our FPGA implementation delivers 1.29$\times$ end-to-end speedup, 1.64$\times$ higher decode throughput, and 3.14$\times$ better energy efficiency on AMD Alveo U280, and an estimated 4.71$\times$, 6.55$\times$, and 4.13$\times$ gains, respectively, on Versal V80.
    \item Leveraging FlexLLM, we build a stage-customized hybrid accelerator for Llama-3.2~1B with a hardware-aware W4A4KV8 SpinQuant~\cite{liu2024spinquant} scheme, reducing WikiText-2 PPL from 13.30 (original SpinQuant) to 12.68. Compared to the BF16 baseline on an NVIDIA A100 GPU, our FPGA implementation delivers 1.29$\times$ end-to-end speedup, 1.64$\times$ higher decode throughput, and 3.14$\times$ better energy efficiency on AMD Alveo U280, and an estimated 4.71$\times$, 6.55$\times$, and 4.13$\times$ gains, respectively, on Versal V80.
    \item We further implement a Hierarchical Memory Transformer (HMT)~\cite{he2024hmt} plug-in on top of FlexLLM for long-context processing, reducing prefill latency by up to 23.23$\times$ and extending the effective context window by over 64$\times$ with less than 7.5\% resource and 0.6\% latency overhead, demonstrating FlexLLM’s scalability to long-context acceleration.
\end{itemize}

% \textbf{Challenge 3: High manual effort and slow iteration for model-specific accelerators.}  
% As LLM architectures evolve rapidly, developing model-specific accelerators remains highly labor-intensive. RTL-based implementations can exceed 100K lines of code and impose a steep learning curve, making it difficult to iterate at the pace of LLM innovation. Recent HLS frameworks such as Allo~\cite{chen2024allo} and StreamTensor~\cite{ye2025streamtensor} reduce RTL burden via MLIR-based mapping and automated compilation, but they largely target spatial architectures and offer limited support for hybrid designs that optimize prefill and decode specifically.

% To close this productivity gap, \textbf{\textit{a high-level programming framework that enables flexible architecture design is necessary for the rapid growth of LLM-oriented accelerators.}} Without a composable framework, accelerator development cycles (often 1--2 years) cannot match the roughly six-month cadence of new LLM generations. What is needed is a framework that exposes the key architectural degrees of freedom—especially stage-customized hybrid design and advanced quantization—while remaining accessible to both ML and accelerator developers.

% Motivated by these challenges, we develop \textbf{FlexLLM}, a composable HLS library built on TAPA~\cite{guo2023tapa}. FlexLLM provides parameterized templates to rapidly construct hybrid accelerators and, to our knowledge, is the first to explicitly reject the unified-design paradigm by enabling stage-customized hybrid architectures, where prefill and decode use different mixes of spatial parallelism and temporal reuse to match their distinct bottlenecks. We will open-source FlexLLM upon publication.

% Using FlexLLM, we build a series of hybrid accelerators for recent LLM workloads with fewer than 1K lines of C++ code. Our design integrates Llama-3.2~1B with state-of-the-art techniques including SpinQuant~\cite{liu2024spinquant} for low-bit quantization and the Hierarchical Memory Transformer (HMT)~\cite{he2024hmt} plug-in for efficient long-context processing. FlexLLM enables new algorithmic innovations to be incorporated into accelerator pipelines quickly, substantially reducing development effort while maintaining high performance and architectural flexibility.

% \noindent\textbf{In summary, our main contributions include:}
% \begin{itemize}[wide=0pt]
%     \item[\textbf{(a)}] FlexLLM: a composable HLS library enabling rapid construction of hybrid LLM accelerators via parameterized templates, substantially reducing design complexity and code volume.
%     \item[\textbf{(b)}] Comprehensive quantization/dequantization support, enabling accurate and efficient low-bit deployment on custom accelerators.
%     \item[\textbf{(c)}] A stage-customized hybrid accelerator for Llama-3.2~1B with a hardware-aware, co-designed W4A4KV8 SpinQuant scheme, achieving 1.29$\times$ end-to-end speedup, 1.64$\times$ higher decode throughput, and 3.14$\times$ better energy efficiency on AMD Alveo U280 versus an NVIDIA A100 80G BF16 baseline (estimated 4.71$\times$/6.55$\times$/4.13$\times$ on Versal V80).
%     \item[\textbf{(d)}] An HMT plug-in integrated via FlexLLM, reducing prefill latency by up to 23.23$\times$ and extending the effective context window by $>64\times$ with $<7.5\%$ resource and $<0.6\%$ latency overhead.
% \end{itemize}


\section{Background}

\subsection{Characterization of LLM Inference Stages}
\label{ssec:llm-inference-stages}

% The Transformer architecture~\cite{vaswani2017attention} forms the computational backbone of most modern LLMs. Each Transformer layer primarily consists of a multi-head self-attention (MHA) block and a feed-forward network (FFN), both of which are followed by normalization and residual connections, shown as Fig.~\ref{fig:arch_comparison}(a). 
% % Owing to this highly regular structure, LLMs exhibit a smaller architectural design space than general ML models, making a kernel-based hardware library a natural and efficient choice for accelerator design.

% During inference, LLMs operate in a decoder-only fashion, where the computation can be divided into two distinct stages: the \textbf{prefill} stage and the \textbf{decode} stage. In the prefill stage, the model processes the entire input prompt in parallel, computing all layer activations for all tokens simultaneously. This stage exhibits high arithmetic intensity and abundant inter-token parallelism, making it largely compute-bound. In contrast, the decode stage generates tokens autoregressively—typically one token at a time—while reusing the stored Key–Value (KV) cache from previous steps to maintain context. As a result, computation in this stage becomes dominated by sequential dependencies and KV memory access, shifting the performance bottleneck from compute throughput to memory bandwidth. 

% To illustrate this difference, we evaluate the BF16 Llama-3.2 1B model~\cite{hf_llama_3_2_1b_access} running on an NVIDIA A100 GPU with vLLM~\cite{kwon2023vllm}, process and generate 1K token in prefill and decode respectively, and exhibit the utilization of compute and memory bandwidth during each stage in  Figure~\ref{fig:prefill_decode_compute_bandwidth}. During the prefill stage, the A100 GPU exhibits a clear speed advantage due to its enormous compute throughput and highly optimized parallel kernels, enabling ultra-fast ingestion of large prompts. However, in the decode stage, the performance bottleneck shifts—compute units become underutilized, and memory bandwidth emerges as the dominant constraint. This comparison highlights the contrast between their resource characteristics.

% Table~\ref{tab:hardware_comparison} summarizes the key hardware characteristics of representative accelerator platforms in comparison~\cite{amd_u280_ds, amd_versal_v80_ds} to the A100 GPU~\cite{nvidia_a100_80gb_pcie_ds}. These results suggest that when effective quantization techniques are applied to reduce activation and weight precision (thereby alleviating memory traffic) and coupled with an accelerator architecture optimized for both computation efficiency and memory reuse, domain-specific accelerators can achieve competitive—and in some cases superior—performance~\cite{chen2024understanding}. This advantage becomes particularly evident in decode-intensive, long-sequence generation workloads such as code synthesis~\cite{joel2024survey} and creative text generation~\cite{gomez2023confederacy}, where token generation latency and energy efficiency are critical. In summary, while GPUs excel in the prefill stage, accelerators co-designed with quantization and hybrid architectural optimizations can match or surpass GPU performance in decode throughput and energy efficiency for inference scenarios.

The Transformer architecture~\cite{vaswani2017attention} forms the backbone of most modern LLMs. Each Transformer layer primarily consists of a multi-head self-attention (MHA) block and a feed-forward network (FFN), followed by normalization and residual connections, as shown in Fig.~\ref{fig:arch_comparison}(a). This highly regular structure makes a kernel-based hardware library a natural and efficient choice for accelerator design.

During inference, decoder-only LLMs are typically executed in two stages: prefill and decode. In prefill, the model processes the entire prompt in parallel and computes activations for all tokens, exhibiting high arithmetic intensity and abundant inter-token parallelism, and is thus largely compute-bound. In decode, the model generates tokens autoregressively while reusing the Key--Value (KV) cache; computation becomes dominated by sequential dependencies and KV access, shifting the bottleneck from compute throughput to memory bandwidth.

To quantify this stage divergence, we run BF16 Llama-3.2~1B~\cite{hf_llama_3_2_1b_access} on an NVIDIA A100 with vLLM~\cite{kwon2023vllm} and measure compute and memory-bandwidth utilization when processing and generating 1K tokens in prefill and decode, respectively (Fig.~\ref{fig:prefill_decode_compute_bandwidth}). The GPU achieves high utilization and strong performance in prefill due to massive parallel compute and highly optimized kernels. In contrast, during decode, compute utilization drops sharply and memory bandwidth becomes the dominant limiter, confirming that the two stages stress fundamentally different resources.

Table~\ref{tab:hardware_comparison} compares representative FPGA platforms~\cite{amd_u280_ds, amd_versal_v80_ds} with the A100 GPU~\cite{nvidia_a100_80gb_pcie_ds}. Notably, the decode stage on GPU often fails to fully utilize its peak bandwidth and compute, whereas modern FPGAs provide ample peak compute and HBM bandwidth to meet decode demands when memory traffic is reduced by quantization. Consequently, with effective low-bit quantization and an architecture co-designed for memory reuse and pipeline efficiency, domain-specific accelerators can be highly competitive—especially for decode-intensive, long-sequence generation workloads where token latency and energy efficiency are critical~\cite{chen2024understanding, joel2024survey, gomez2023confederacy}. 
% In summary, GPUs excel in prefill, while accelerators co-designed with quantization and hybrid architectural optimizations can match or surpass GPUs in decode throughput and energy efficiency.



\begin{table}[t]
\centering
\caption{Comparison of key hardware metrics across NVIDIA A100 GPU, AMD Alveo U280 FPGA, and Versal V80 FPGA.}
\resizebox{1.0\linewidth}{!}{%
    \begin{tabular}{l|c|c|c}
    \hline
    \textbf{Metric}     & \textbf{U280} & \textbf{V80}  & \textbf{A100 80GB PCle} \\
    \hline
    Technology Node     & TSMC 16nm         & TSMC 7nm              & TSMC 7nm   \\
    Peak Compute        & 8 FP32 TFLOPS     & 58 FP32 TFLOPS      & 312 FP32 TFLOPS   \\
    Peak HBM BW         & 460 GB/s          & 820 GB/s              & 1,935 GB/s    \\
    HBM Capacity        & 8 GB              & 32 GB                 & 80 GB         \\
    Peak Power          & 75 W             & 190 W                 & 300 W         \\
    \hline
    \end{tabular}%
}
\label{tab:hardware_comparison}
\vspace{-0.2cm}
\end{table}

\begin{figure}[t]
    \centering
    \includegraphics[width=0.75\linewidth]{Figures/Prefill_vs_Decode_Compute_Bandwidth_Legend36.pdf}
    \vspace{-0.3cm}
    \caption{Profiling of compute throughput and memory bandwidth utilization during the prefill and decode stages of the BF16 Llama-3.2 1B model on an NVIDIA A100 GPU.}
    \label{fig:prefill_decode_compute_bandwidth}
    \vspace{-0.4cm}
\end{figure}



\begin{table*}[t]
\centering
\caption{Comparison of related accelerators and frameworks.}
\label{tab:framework-comparison}
\vspace{-0.2cm}
\resizebox{\textwidth}{!}{
    \begin{tabular}{lcccccc}
    \hline
    Accelerator Framework &
    Contribution Type &
    Arch. Type &
    Extensibility &
    stage-customized kernel Opt. &
    Support Precision &
    Support Quant. Type \\
    \hline
    FlightLLM~\cite{zeng2024flightllm} & Accelerator & Temporal & Not extensible & \xmark & FP32, INT1$\sim$8 & Static \\
    FINN~\cite{blott2018finn}, hls4ml~\cite{fastml_hls4ml} & Compiler Framework & Spatial & Automatic Mapping & \xmark & FP32, INT1$\sim$8 & Static \\
    Allo~\cite{chen2024allo} & Domain-Specific Language & Spatial & Require Manual Design & \xmark & FP32, INT4/8 & Static \\
    StreamTensor~\cite{ye2025streamtensor} & Compiler Framework & Spatial & Automatic Optimization & \xmark & FP32, INT8 & Static \\
    SSR~\cite{zhuang2024ssr}, InTAR~\cite{he2025intar} & Accelerator & Hybrid & Require Manual Design & \xmark & FP32, INT8 & Static \\
    
    \hline
    FlexLLM (Ours) & HLS Kernel Library & Hybrid & Templated Kernel Composition & \cmark & FP16/32, INT1$\sim$8 & Static/Dynamic \\
    \hline
    \end{tabular}%
}
\vspace{-0.4cm}
\end{table*}






\subsection{Hardware-efficient LLM Quantization}
% Extensive research has been devoted to quantization for efficient LLM deployment on domain-specific accelerators. The general form of $N$-bit quantization can be expressed as:
% \[
% X_q = s \cdot \text{round}\left(\frac{X - b}{s}\right) + b,
% \]
% where $X$ denotes the original floating-point tensor, $X_q$ represents the quantized tensor, and $s$ and $b$ are the scaling and offset parameters, respectively. In symmetric quantization, $s = \frac{\max(|X|)}{2^{N-1}-1}$ and $b = 0$, while in asymmetric quantization, $s = \frac{\max(X) - \min(X)}{2^{N}-1}$ and $b = \min(X)$. Quantization is referred to as dynamic when $s$ and $b$ are measured at runtime, and static when they are precomputed statistically during offline calibration.

% While static quantization is naturally compatible with pre-trained weights during inference, it is also more hardware-friendly for activations since it avoids complex floating-point operations. Consequently, it has been widely adopted in recent accelerator designs~\cite{zeng2024flightllm,chen2024understanding}. However, its effectiveness deteriorates notably when scaling down to lower bit-widths such as INT4, even under advanced schemes like SmoothQuant~\cite{xiao2023smoothquant}. In contrast, finer-grained dynamic quantization—when combined with outlier-handling techniques such as activation rotation and Fast Hadamard Transform (FHT) introduced in SpinQuant~\cite{liu2024spinquant}—offers a more practical and accurate solution for aggressive low-bit quantization in LLM deployment. It is important to note that the per-token \textit{max} and \textit{min} required in dynamic quantization not only incur additional computational overhead but also introduce token-wise dependency barriers similar to those in LayerNorm. These dependencies complicate hardware pipeline scheduling, thereby demanding greater flexibility in both architectural design and runtime optimization.

Quantization is essential for efficient LLM deployment on domain-specific accelerators. A general $N$-bit integer quantization can be written as
\[
X_q = s \cdot \text{round}\left(\frac{X - b}{s}\right) + b,
\]
where $X$ is the floating-point (FP) tensor, $X_q$ is the quantized integer tensor, and $s$ and $b$ are the scale and zero offset. Symmetric quantization uses $s = \frac{\max(|X|)}{2^{N-1}-1}$ and $b = 0$, while asymmetric quantization uses $s = \frac{\max(X) - \min(X)}{2^{N}-1}$ and $b = \min(X)$. Quantization is \emph{static} when $s,b$ are precomputed offline, and \emph{dynamic} when they are measured at runtime.

Static quantization is naturally compatible with pre-trained weights and is hardware-friendly for activations since it avoids runtime calibration, and thus has been widely adopted in recent accelerators~\cite{zeng2024flightllm,chen2024understanding}. However, its accuracy often degrades sharply at INT4, even with advanced schemes such as SmoothQuant~\cite{xiao2023smoothquant}. In contrast, fine-grained dynamic quantization introduces extra overhead and token-wise dependency barriers, but becomes a practical route to accurate aggressive low-bit deployment when combined with outlier-handling techniques, such as activation rotation and Fast Hadamard Transform (FHT) introduced in SpinQuant~\cite{liu2024spinquant}.


\subsection{Related Accelerators and Frameworks}

Table~\ref{tab:framework-comparison} summarizes representative Transformer-oriented accelerators and frameworks.

FlightLLM~\cite{zeng2024flightllm} adopt a \textbf{temporal} design that sequentially executes layers on shared engines to maximize reuse. While utilization can be high, limited on-chip buffering incurs frequent off-chip accesses, and the monolithic engine limits support for heterogeneous kernels/precisions and rapid model customization, increasing latency, energy, and redesign effort. 

In contrast, FINN~\cite{blott2018finn}, hls4ml~\cite{fastml_hls4ml}, Allo~\cite{chen2024allo} and StreamTensor~\cite{ye2025streamtensor} embrace a fully \textbf{spatial-dataflow} paradigm, mapping modules to dedicated hardware and streaming data on-chip. They provide useful abstractions or compilation support for customization, but decode performance remains limited by pipeline stalls under intrinsic dependencies and unbalanced kernel latencies, where throughput becomes hard to sustain.

Hybrid architectures aim to combine the strengths of temporal and spatial designs via flexible scheduling and module-level customization. For example, SSR~\cite{zhuang2024ssr} and InTAR~\cite{he2025intar} instantiate a fixed set of linear and non-linear engines, enabling inter-module streaming and intra-engine kernel switching. However, existing hybrids largely retain a \emph{unified} architecture across inference stages; fixed engine configurations restrict deeper optimization across diverse LLM models and the sharply different bottlenecks of prefill and decode.

To fully realize the potential of hybrid LLM acceleration, \textbf{stage-customized} design is necessary, where prefill and decode adopt different module configurations and scheduling policies to match their distinct hardware characteristics (Fig.~\ref{fig:arch_comparison}(f)(g)). In the next section, we describe how FlexLLM achieve this goal through stage-oriented module optimization and rapid hybrid construction, while integrating broad quantization support to build practical high-performance LLM accelerators with a model-to-silicon turnaround of less than two months.



% \subsection{Motivation for FlexLLM Design}

% % Building upon the above reviews, we identify three design motivations that guide the development of FlexLLM.


% % \textbf{stage-customized module optimization.}  
% % As discussed in Section~\ref{fig:prefill_decode_compute_bandwidth}, the prefill and decode stages of LLM inference exhibit fundamentally different computational and memory characteristics. Therefore, accelerator modules must be designed and optimized separately for each stage to achieve balanced performance and efficient resource utilization across the entire inference process.

% % \textbf{Comprehensive quantization support for practical accelerator deployment.}  
% % Quantization is critical for reducing computation and memory overhead in domain-specific accelerators, yet existing frameworks often support only limited quantization schemes or precision levels. A practical and scalable framework must flexibly support a wide range of quantization strategies, while preserving model accuracy under aggressive low-bit configurations.

% % \textbf{Flexible hybrid architecture exploration with fast implementation flow.} 
% % A hybrid design space with composable configuration can better exploits the trade-offs between temporal and spatial architectures, each offering distinct strengths in utilization and parallelism. Moreover, long RTL development and physical design cycles remain major bottlenecks for deploying new models. An efficient HLS-based design flow is therefore essential to shorten the model-to-silicon timeline from several months to weeks.

% Building on the above analysis, we identify three key motivations that guide the development of FlexLLM.

% \textbf{stage-customized module optimization.}  
% The prefill and decode stages of LLM inference exhibit fundamentally different compute and memory behaviors. Achieving balanced performance and efficient resource utilization therefore requires modules to be tailored and optimized for each stage individually.

% \textbf{Comprehensive quantization support.}  
% Quantization is vital for reducing computation and memory overheads, yet existing frameworks support limited precision and quantization schemes. A practical framework should flexibly support diverse quantization strategies while maintaining accuracy under aggressive low-bit settings.

% \textbf{Flexible hybrid architecture with fast implementation.}  
% A composable hybrid design space can better exploit the trade-offs between temporal and spatial architectures, each excelling in utilization or parallelism. Moreover, efficient HLS-based flows are essential to overcome long RTL and physical design cycles, reducing the model-to-silicon timeline from months to weeks.


% Motivated by these needs, we propose \textbf{FlexLLM}, a composable HLS-based hybrid library that supports flexible module configuration, stage-customized optimization, and diverse quantization schemes. Built on the TAPA HLS framework, FlexLLM offers a highly parameterized design flow and seamless compatibility with AutoBridge~\cite{guo2021autobridge}, enabling rapid design construction and automatic placement-and-routing optimization. Together, these capabilities achieve a model-to-silicon turnaround time of less than two months.

% \section{FlexLLM Framework}
% \subsection{stage-customized Module Templates}

% While the prefill stage exposes abundant parallelism across tokens, the decode stage retains parallelism only within a single token’s computation due to autoregressive data dependencies. Based on these distinct characteristics, we apply stage-specific customization strategies to maximize hardware efficiency.

% \textbf{Prefill Stage Modules.}  
% During prefill, activations from multiple tokens are packed and processed concurrently by hardware modules for both linear and non-linear layers, as shown in Figure~\ref{fig:moudle}(a). We define this configurable inter-token parallelism as \textit{\textbf{token\_parallelism (TP)}}. Since linear layers dominate the overall computational workload, we further introduce a configurable \textit{\textbf{weight\_parallelism (WP)}} that enables multiple weight channels to be fetched and streamed simultaneously from off-chip memory.

% Each linear layer module therefore forms a two-dimensional systolic array of \textit{TP}$\times$\textit{WP} processing elements (PEs). These PEs are optimized for precision-specific arithmetic operations, supporting packed integer computation and achieving an pipeline initiation interval (II) of one cycle on the target platform, improving both throughput and energy efficiency.

% Given a prompt of sequence length $l_{p}$, input dimension $d_{in}$, and output dimension $d_{out}$, the theoretical latency of a linear layer in the prefill stage can be expressed as:
% \begin{equation}
% T_{linear}^{p} = \frac{l_{p} \times d_{in} \times d_{out}}{\textit{TP} \times \textit{WP}}
% \label{eq:linear_latency_prefill}
% \end{equation}
% The off-chip memory bandwidth consumption is:
% \begin{equation}
% BW_{linear}^{p} = B_W \times \textit{WP} \times F,
% \label{eq:linear_bandwidth}
% \end{equation}
% where $B_W$ denotes the number of bytes per weight element and $F$ is the achieved accelerator frequency. This model allows users to customize the \textit{TP} and \textit{WP} parameters of each module to balance throughput across layers, manage memory bandwidth utilization, and maintain a well-pipelined execution flow.

% \textbf{Decode Stage Modules.}  
% During decoding, we exploit intra-token parallelism by dividing the calculation of each output hidden vector into multiple chunks processed in parallel and subsequently aggregated through on-chip buffers or reduction trees. We define this form of parallelism as \textit{\textbf{block\_parallelism (BP)}}. Each linear layer also supports similar \textit{\textbf{WP}} to further enhance throughput. Inside the module, vector blocks are processed by \textit{BP} one-dimensional systolic arrays, each composed of $\textit{WP}/\textit{BP}$ PEs, shown as Figure~\ref{fig:moudle} (b). Accordingly, for an output of sequence length $l_{d}$, and the same $d_{in}$ and $d_{out}$, the theoretical latency of a linear layer in the decode stage is:
% \begin{equation}
% T_{linear}^{d} = \frac{l_{d} \times d_{in} \times d_{out}}{\textit{WP}} 
% \label{eq:linear_latency_decode}
% \end{equation}
% while the off-chip bandwidth requirement follows the same formulation as Equation~\ref{eq:linear_bandwidth}. It is noteworthy that, in the decode stage, PEs do not share weight channels as in the prefill stage. Consequently, under the same hardware resource budget, the linear layer module can adopt a higher \textit{WP}, leading to increased memory bandwidth demand during decoding.

% \begin{figure}[t]
%     \centering
%     \includegraphics[width=0.9\linewidth]{Figures/Module_Example.pdf}
%     \vspace{-0.3cm}
%     \caption{stage-customized module design for prefill(a), decode (b), and quantization module integration (c).}
%     \label{fig:moudle}
%     \vspace{-0.4cm}
% \end{figure}

% \textbf{Quantization Modules.}  
% Quantization modules adopt the same parallelism configuration as the non-linear modules. As shown in Figure~\ref{fig:moudle}(c), during quantization, the quantization engine scales floating-point inputs to integers of specific bit widths using scaling factors and zero points, which are either preloaded from memory (static) or computed online (dynamic). After kernel computation, the dequantization engine restores outputs to floating-point form using the same scaling parameters and auxiliary data such as per-channel weight scales and sums buffered on-chip. 
% In general, our framework supports static/dynamic, symmetric/asymmetric quantization types, with per-tensor, per-token, and per-channel granularities, and also includes outlier-handling modules such as rotation and FHT~\cite{liu2024spinquant}.

% \textbf{Library Implementation.}  
% Following these design principles, we develop a comprehensive module library named FlexLLM, built on top of the TAPA HLS framework and comprising over 10K lines of parameterized code. FlexLLM integrates a rich set of compute kernels commonly used in LLMs, along with versatile quantization and dequantization modules supporting a wide spectrum of quantization schemes. It also includes auxiliary components for on-chip streaming, buffering, and memory access management. Table~\ref{tab:module_template} summarizes the main modules and their configurable parameters. These highly templated modules enable extensive customization and composability for hybrid architecture exploration while maintaining ease of use for non-expert developers, effectively reducing the model-to-silicon design cycle from months to weeks.



\section{FlexLLM Framework}
\subsection{stage-customized Module Templates}

As discussed in ~\ref{ssec:llm-inference-stages}, prefill and decode stress fundamentally different resources, so a single unified architecture inevitably compromises one stage. FlexLLM addresses this by providing \textbf{stage-customized module templates} that expose the right parallelism knobs for each stage, enabling rapid construction of hybrid accelerators with balanced pipelines.

\textbf{Prefill-stage modules:}  
Prefill exposes abundant inter-token parallelism. FlexLLM packs activations from multiple tokens and processes them concurrently for both linear and non-linear layers (Fig.~\ref{fig:moudle}(a)). We define this inter-token parallelism as \textit{\textbf{token\_parallelism (TP)}}. Since linear layers dominate compute, we further introduce \textit{\textbf{weight\_parallelism (WP)}} to fetch and stream multiple weight channels from off-chip memory in parallel. Each prefill linear module implements a 2D systolic array of $\textit{TP}\times\textit{WP}$ processing elements (PEs), optimized with an initiation interval (II) of one cycle for supported precisions. 

Given a prompt of length $l_{p}$, the theoretical prefill latency of a linear layer with input dimension $d_{in}$ and output dimension $d_{out}$ can be expressed as
\begin{equation}
T^{p}_{\text{linear}}=\frac{l_p\, d_{in}\, d_{out}}{\textit{TP}\cdot\textit{WP}},
\label{eq:linear_latency_prefill}
\end{equation}
and the corresponding off-chip bandwidth demand is
\begin{equation}
BW^{p}_{\text{linear}} = B_W \cdot \textit{WP} \cdot F,
\label{eq:linear_bandwidth}
\end{equation}
where $B_W$ is bytes per weight element and $F$ is the operating frequency. These knobs allow users to balance layer throughput while respecting memory-bandwidth constraints.

\textbf{Decode-stage modules: }  
Decode is constrained by autoregressive dependencies, leaving parallelism primarily within a single token. FlexLLM therefore exploits intra-token parallelism by partitioning each output hidden vector into blocks computed in parallel and reduced on-chip. We denote this as \textit{\textbf{block\_parallelism (BP)}}. Each decode linear module also supports \textit{\textbf{WP}}, implemented as $\textit{BP}$ sets of 1D systolic arrays with $\textit{WP}/\textit{BP}$ PEs each (Fig.~\ref{fig:moudle}(b)). For decode length $l_d$, the idea latency of a linear layer with the same $d_{in}$ and $d_{out}$ is
\begin{equation}
T^{d}_{\text{linear}}=\frac{l_d\, d_{in}\, d_{out}}{\textit{WP}}.
\label{eq:linear_latency_decode}
\end{equation}
The off-chip bandwidth follows Eq.~\ref{eq:linear_bandwidth}. Compared to prefill, decode PEs do not share weights across tokens, so the same resource budget can often support a higher $\textit{WP}$, increasing bandwidth demand and motivating careful stage-specific tuning.

\begin{figure}[t]
    \centering
    \includegraphics[width=0.9\linewidth]{Figures/Module_Example.pdf}
    \vspace{-0.3cm}
    \caption{stage-customized module design for (a) prefill, (b) decode, and (c) quantization module integration.}
    \label{fig:moudle}
    \vspace{-0.4cm}
\end{figure}

\textbf{Quantization modules: }
Quantization modules use the same configurable parallelism (\textit{TP} for prefll and \textit{BP} for decode) as the non-linear modules. As shown in Figure~\ref{fig:moudle}(c), the quantizer converts FP inputs to low-bit integers using scales and zero offsets, which are either preloaded (static) or computed online (dynamic). After kernel computation, the dequantizer reconstructs FP outputs using the same scales and offsets, along with auxiliary data (e.g., per-channel weight scales and sums) buffered on-chip. In general, our framework supports static/dynamic and symmetric/asymmetric quantization with per-tensor, per-token, and per-channel granularities, and includes outlier-handling modules such as rotation and FHT.


\textbf{Library components: }
Following these principles, FlexLLM provides a comprehensive module library built on TAPA, comprising over 10K lines of highly parameterized code. It includes core LLM kernels, quantization/dequantization modules, and auxiliary components for on-chip streaming, buffering, and memory access management. Table~\ref{tab:module_template} summarizes the key modules and configurable parameters. These templates enable rapid, composable hybrid architecture exploration while remaining accessible to non-expert developers, reducing model-to-silicon turnaround from months to weeks.

\begin{table*}[t]
\centering
\scriptsize
\caption{Overview of primary module templates, configurable parameters, and Module interfaces provided in FlexLLM.}
\vspace{-0.2cm}
\renewcommand{\arraystretch}{1.1}
\setlength{\tabcolsep}{3.4pt}
\begin{tabular}{m{1cm}|m{1.8cm}|m{9.4cm}|m{4.8cm}}
\hline
 & \textbf{Module Template} & \textbf{Module Configurable Parameter} & \textbf{Module Interface} \\
\hline
\multirow{3}{=}{\raggedright\arraybackslash \textbf{Kernel Library}} 
& Linear Layer & 
\texttt{dtype, token\_parallelism(prefill), block\_parallel(decode), head\_parallel(MHA), weight\_parallelism, head\_num(MHA), max\_in\_dim, max\_out\_dim, max\_seq\_len} &
\texttt{in\_stream, w\_stream, out\_stream, in\_dim, out\_dim, seq\_len} \\
\cline{2-4}
& Non-Linear Layer (RoPE, Softmax, LayerNorm, ...) &
\texttt{dtype, token\_parallelism(prefill), block\_parallel(decode), head\_parallel(MHA), head\_num(MHA), max\_io\_dim, max\_seq\_len} &
\texttt{in\_stream, out\_stream, io\_dim, seq\_len} \\
\hline
\multirow{3}{=}{\raggedright\arraybackslash \textbf{Quant Library}}
& Static/Dynamic Quant Layer &
\texttt{in\_dtype, in\_quant\_bit, token\_parallelism(prefill), block\_parallel(decode), head\_parallel(MHA), in\_quant\_type(sym/asym), in\_quant\_granularity(per-tensor/token), head\_num(MHA), max\_io\_dim, max\_seq\_len} &
\texttt{in\_stream, in\_scale\_stream, in\_zero\_stream, quant\_in\_stream, io\_dim, seq\_len} \\
\cline{2-4}
& Static/Dynamic Dequant Layer &
\texttt{in\_quant\_bit, w\_quant\_bit, out\_dtype, token\_parallelism(prefill), block\_parallel(decode), head\_parallel(MHA), in\_quant\_type(sym/asym), in\_quant\_granularity(per-tensor/token), w\_quant\_type(sym/asym), w\_quant\_granularity(per-tensor/channel), head\_num(MHA), max\_io\_dim, max\_seq\_len} &
\texttt{quant\_out\_stream, in\_scale\_stream, in\_zero\_stream, w\_scale\_stream, w\_col\_sum\_stream, out\_stream, io\_dim, seq\_len} \\
\hline
\end{tabular}
% \begin{minipage}{1.0\linewidth}
% \scriptsize
% \vspace{3pt}
% \hspace{15pt}$\wedge$ serves as input/output stream when static/dynamic scaling is used. \quad
% $*$ enabled when asymmetric quantization is applied to activations.
% \end{minipage}
\label{tab:module_template}
\vspace{-0.4cm}
\end{table*}


\subsection{Hybrid Accelerator Construction \& Optimization}

Built upon the FlexLLM and TAPA HLS flows, users can efficiently explore hybrid architectures for LLM accelerators and quickly transition from design exploration to on-board implementation. Figure~\ref{fig:code} illustrates a simplified example that composes several modules in a hybrid style.

In the temporal-reuse part, the same templated module is instantiated and invoked multiple times within a single function to sequentially process similar operations. In this example, the prefill module of the linear layer and RoPE is reused for both \textit{Key} and \textit{Query} computations. This approach maximizes hardware utilization by reusing computation resources across similar tasks while minimizing redundant module instantiations.

In the spatial dataflow part, each instantiated module is explicitly invoked at the top level and connected through on-chip FIFOs to enable parallel execution and streaming communication. This module-based composition exploits inter-module pipelining and on-chip dataflow to achieve high throughput and overlap between computation stages.

The combination of temporal and spatial styles can effectively balance performance and resource efficiency in hybrid accelerator design. We illustrate a detailed methodology for stage-customized hybrid design through a case study in Sec.~\ref{case_study_1}. After composing an accelerator with FlexLLM, users can leverage its seamless compatibility with AutoBridge~\cite{guo2021autobridge} to optimize placement and routing (P\&R), enabling parallel exploration of design candidates to achieve higher frequency and performance. This end-to-end toolchain accelerates iteration and supports efficient on-board deployment.

\begin{figure}[t]
    \centering
    \includegraphics[width=0.8\linewidth]{Figures/code_example.png.pdf}
    \vspace{-0.4cm}
    \caption{Example code illustrating hybrid architecture construction combining temporal reuse and spatial dataflow with FlexLLM.}
    \label{fig:code}
    \vspace{-0.8cm}
\end{figure}

\begin{figure*}[t]
    \centering
    \includegraphics[width=0.9\linewidth]{Figures/SpinQuant_HMT_Design.pdf}
    \vspace{-0.2cm}
    \caption{Hybrid architecture design for (a) prefill \& (b) decode, and  (c) the integration diagram of HMT plug-in.}
    \label{fig:Implementation}
    \vspace{-0.4cm}
\end{figure*}

\section{Case Study 1: Llama 3.2-1B + SpinQuant}
\label{case_study_1}

% In this section, we demonstrate how FlexLLM is utilized to construct a comprehensive high-performance accelerator system for the Llama-3.2 1B model. Specifically, we first revise SpinQuant to obtain a more hardware-efficient representation while minimizing accuracy degradation. We then compose stage-customized accelerator architectures using FlexLLM's composable modules. To further mitigate the quadratic computation and memory overhead in long-context processing, we integrate the Hierarchical Memory Transformer (HMT) plug-in into the backbone model, transforming both compute and memory complexity into a linear relation with respect to prompt length. Table~\ref{tab:module_summary} summarizes the major module usages in each component, along with their corresponding lines of code and estimated development workload.

In this section, we demonstrate how FlexLLM is utilized to construct a high-performance accelerator with stage-customized architectures for the Llama-3.2 1B model. Specifically, we first refine the SpinQuant framework to obtain a more hardware-efficient representation while minimizing accuracy degradation. We then build customized architectures for prefill and decode with composable modules from FlexLLM. Table~\ref{tab:module_summary} summarizes the major module components, along with their corresponding lines of code and estimated development workload.


\subsection{Hardware-Efficient Model Quantization}

% The original SpinQuant framework enhances quantization accuracy by removing activation outliers through learned rotations\cite{liu2024spinquant}. While most rotations can be merged into the quantized weights or efficiently implemented using FHT, the remaining boundary rotations still require FP computation comparable to a full linear layer—an undesirable overhead for targeted FPGA platforms. In addition, the baseline implementation does not quantize the query vectors in the MHA and the final vocabulary projection (lm\_head) layer, retaining them in FP that would severely reduce hardware efficiency.

% To enable a more hardware-efficient model, we apply the following optimizations while maintaining accuracy as much as possible:

%   \textbf{Removing boundary rotations.}  
%   We eliminate the first and last rotations in each Transformer block by folding their effects into neighboring weights during fine-tuning. This modification preserves the outlier-removal benefits while completely removing floating-point boundary operations.

%   \textbf{Quantized attention at higher bit-width.}  
%   Because MHA is highly sensitive to precision, we use INT8 for attention while keeping other linear layers at INT4. For hardware simplicity, MHA adopts \emph{static symmetric per-tensor} quantization, whereas the other linear layers employ \emph{dynamic asymmetric per-token} quantization to retain accuracy. The online Hadamard rotation within MHA is removed as it contributes little to accuracy under this configuration.

%   \textbf{Integer vocabulary projection.}  
%   The lm\_head layer is quantized to INT4, consistent with other linear layers, reducing resource consumption and improving decoding throughput.

% Table~\ref{tab:ablation} summarizes the ablation study on WikiText-2. Starting from the BF16 (PPL = 8.94), the original SpinQuant configuration with INT4 linear layers and BF16–INT4 attention increases PPL to 13.30. Elevating attention precision to INT8 (dynamic) recovers considerable lost accuracy. Switching to static INT8 slightly relaxes accuracy but simplifies hardware implementation. Finally, quantizing the lm\_head layer introduces only a minor PPL increase while enabling a fully integer pipeline across all linear layers.

SpinQuant improves low-bit accuracy by mitigating activation outliers via learned rotations~\cite{liu2024spinquant}. While many rotations can be folded into weights or implemented efficiently with FHT, the remaining boundary rotations still require FP compute comparable to a full linear layer, which is costly on FPGAs. Moreover, the baseline SpinQuant setup keeps the MHA query path and the final vocabulary projection (\texttt{lm\_head}) in FP, leaving substantial hardware efficiency unrealized.

To obtain a practical model with minimal accuracy loss, we make the following hardware-oriented refinements:

\textbf{Remove boundary rotations.}  
We eliminate the first and last rotations in each decoder layer by absorbing them into adjacent weights during fine-tuning. This preserves the outlier-mitigation effect while removing FP boundary operations entirely.

\textbf{Quantize attention with higher precision.}  
Since MHA is more precision-sensitive, we quantize attention to INT8 while keeping other linear layers at INT4. For hardware simplicity, MHA uses static symmetric per-tensor quantization, while the remaining linear layers use dynamic asymmetric per-token quantization to retain accuracy. 
% We also remove the online Hadamard rotation in MHA, as it provides negligible benefit under this configuration.

\textbf{Integer vocabulary projection.}  
We further quantize \texttt{lm\_head} to INT4, matching the other linear layers, reducing resource cost and improving decode throughput.

Table~\ref{tab:ablation} reports an ablation on WikiText-2. Starting from BF16 (PPL = 8.94), the original SpinQuant setup (INT4 linear layers with BF16-INT4 attention) yields PPL = 13.30. Raising attention precision to INT8 recovers most of the lost accuracy; using static INT8 slightly relaxes PPL but simplifies hardware. Finally, quantizing \texttt{lm\_head} increases PPL only marginally while enabling a fully integer pipeline across all linear layers.


\subsection{Architecture Design}

% Leveraging the stage-customized modules provided by FlexLLM, we implement separate architectures for the prefill and decode stages of the quantized model to optimize performance and resource utilization for each phase.

% \textbf{Prefill stage.}  
% Unlike a purely spatial design that maps each kernel in a decoder layer to a dedicated hardware module, our prefill implementation introduces an additional temporal-reuse strategy. Specifically, the same linear layer and RoPE modules are shared between Q\_proj and K\_proj, similar for V\_proj and O\_proj, as illustrated in Figure~\ref{fig:Implementation}(a). The \textit{Key} and \textit{Value} tensors are computed in parallel and stored in off-chip HBM, after which the remaining layers execute in a fully pipelined dataflow fashion. The streaming connections between modules maximize pipeline overlap across tokens, while temporal reuse eliminates idleness caused by data dependencies and help form a balanced pipeline in MHA.

% Following Eq.~\ref{eq:linear_latency_prefill}, the theoretical latency bound of the prefill stage can be expressed as:
% \begin{equation} \label{eq:prefill_latency}
%     T_{p} = 
% \frac{Nl_{p}}{\textit{TP}} ( 
% \frac{d_{h}d_{\textit{kv}}}{\textit{WP}_{\textit{kqvo}}} + 
% \max( 
% \frac{d_{h}^2}{\textit{WP}_{\textit{kqvo}}}, 
% \frac{d_{h}l_{p}}{\textit{WP}_{\textit{mha}}}, 
% \frac{d_{h}d_{\textit{ffn}}}{\textit{WP}_{\textit{ffn}}}))
% \end{equation}
% where $N$ denotes the number of decoder layers, and $d$, $d_{kv}$, and $d_{ffn}$ represent the hidden dimension, key/value dimension, and FFN intermediate dimension, respectively. $\textit{WP}_{\textit{kqvo}}$, $\textit{WP}_{\textit{mha}}$, and $\textit{WP}_{\textit{ffn}}$ represent the weight parallelism configurations for different linear modules.

% The corresponding peak memory bandwidth is:
% \begin{equation} \label{eq:prefill_bandwidth}
% BW_{p} = F(B_W^{\textit{int4}}(2\textit{WP}_{\textit{kqvo}} + 3\textit{WP}_{\textit{ffn}}) + B_W^{\textit{int8}} 2\textit{WP}_{\textit{mha}}) 
% \end{equation}
% where $B_W$ represents the bandwidth per weight element. 

% While the overhead of non-linear modules scales primarily with \textit{TP}, that of linear modules scales with both \textit{TP} and \textit{WP}. By tuning these parameters—for example, ensuring 
% $\frac{d_{h}^2}{\textit{WP}_{\textit{kqvo}}} =
% \frac{d_{h}L_p}{\textit{WP}_{\textit{mha}}} = 
% \frac{d_{h}d_{\textit{ffn}}}{\textit{WP}_{\textit{ffn}}}$— 
% and optimizing \textit{TP} and \textit{WP} values via Integer Linear Programming (ILP) under hardware constraints, users can balance compute throughput and memory utilization to achieve an optimized prefill pipeline.

Using FlexLLM's stage-customized modules, we implement separate architectures for prefill and decode to match their distinct bottlenecks and achieve better performance and efficiency.

\textbf{Prefill stage:}  
Prefill benefits from inter-token parallelism and pipelined streaming. Unlike a purely spatial design that instantiates a dedicated module for each kernel, our prefill architecture combines streaming dataflow with selective temporal reuse. Shown as Fig.~\ref{fig:Implementation}(a), we share the same linear and RoPE modules between \texttt{Q\_proj} and \texttt{K\_proj}, and similarly for \texttt{V\_proj} and \texttt{O\_proj}. We compute \textit{Key} and \textit{Value} in parallel and store them in off-chip HBM, after which the remaining kernels execute in a fully pipelined dataflow manner. Streaming links maximize overlap across tokens, while temporal reuse mitigates idleness caused by data dependencies and helps form a balanced pipeline in MHA.

Following Eq.~\ref{eq:linear_latency_prefill}, the prefill latency bound is
\begin{equation} \label{eq:prefill_latency}
    T_{p} = 
\frac{Nl_{p}}{\textit{TP}} ( 
\frac{d_{h}d_{\textit{kv}}}{\textit{WP}_{\textit{kqvo}}} + 
\max( 
\frac{d_{h}^2}{\textit{WP}_{\textit{kqvo}}}, 
\frac{d_{h}l_{p}}{\textit{WP}_{\textit{mha}}}, 
\frac{d_{h}d_{\textit{ffn}}}{\textit{WP}_{\textit{ffn}}}))
\end{equation}
where $N$ is the number of decoder layers, and $d_h$, $d_{\textit{kv}}$, and $d_{\textit{ffn}}$ are the hidden, KV, and FFN intermediate dimensions. $\textit{WP}_{\textit{kqvo}}$, $\textit{WP}_{\textit{mha}}$, and $\textit{WP}_{\textit{ffn}}$ denote the weight-parallelism settings of the corresponding modules. The peak bandwidth demand is
\begin{equation}\label{eq:prefill_bandwidth}
BW_{p}=F\!\left(B_W^{\textit{int4}}\!\left(2\textit{WP}_{\textit{kqvo}}+3\textit{WP}_{\textit{ffn}}\right)+B_W^{\textit{int8}}\cdot 2\textit{WP}_{\textit{mha}}\right)
\end{equation}
where $B_W$ is bytes per weight element and F denotes frequency.

In prefill, non-linear overheads scale mainly with \textit{TP}, while linear modules scale with both \textit{TP} and \textit{WP}. We tune these parameters via Integer Linear Programming (ILP) under hardware constraints (resources and memory bandwidth) to minimize $T_{p}$ and obtain an optimal, well-pipelined design.

\textbf{Decode stage:}  Autoregressive dependencies prevent inter-token parallelism in decode, making spatial overlap across tokens infeasible. To maximize runtime hardware utilization, we temporally reuse modules for all INT4 quantization, linear layers, and non-linear layers, while maintaining dataflow connectivity within MHA and between reused quantization, linear, non-linear modules, shown as Fig.~\ref{fig:Implementation}(b). This design enables intra-token parallelism and inter-head overlay, effectively exploiting available concurrency.

The theoretical latency bound and memory bandwidth of the decode stage are given in Eq.~\ref{eq:decode_latency} and Eq.~\ref{eq:decode_bandwidth}, where $\textit{WP}_{int4}$ denotes the shared \textit{WP} used for the projection, FFN, and \texttt{lm\_head} layers. We similarly tune \textit{BP}s and \textit{WP}s via ILP under hardware constraints to minimize $T{d}$ on the target device.
\begin{equation} \label{eq:decode_latency}
    \begin{split}
        T_{d} &= 
        l_{d}  
        (
        \frac{N(2 d_{h}d_{\textit{kv}} + d_{h}^2 + 3 d_{h}d_{\textit{ffn}}) + dd_{\textit{lm\_head}}}
        {\textit{WP}_{int4}} \\
        &+ \max(
        \frac{Nd_{h}^2}{\textit{WP}_{int4}},
        \frac{Nd_{h} (l_{p} + 0.5 l_{d})}
        {\textit{WP}_{\textit{mha}}}
        )
        )
    \end{split}
\end{equation}
\begin{equation}
BW_{d} = 
F(B_W^{\textit{int4}}\textit{WP}_{\textit{int4}} + 
2B_W^{\textit{int8}}\textit{WP}_{\textit{mha}})
\label{eq:decode_bandwidth}
\end{equation}

% In practical implementation, the \textit{WP} setting in the decode stage is often higher than in prefill, resulting in large modules with substantial signal fan-out. To address this, we divide oversized linear-layer modules into multiple identical submodules, improving floorplanning flexibility and increasing achievable frequency during place-and-route.

In practice, decode often uses larger \textit{WP} setting, creating oversized linear modules with high fan-out. We mitigate this by partitioning them into multiple identical submodules, improving floorplanning flexibility and achievable frequency during P\&R.

\begin{table}[t]
\centering
\scriptsize
\caption{Summary of module usage, code size, and development workload for each design component.}
\vspace{-0.2cm}
\renewcommand{\arraystretch}{1.1}
\setlength{\tabcolsep}{2.5pt}
\begin{tabular}{l|p{3.5cm}|c|c}
\hline
\textbf{Design} & \textbf{Main Module Usage} & \textbf{LoC} & \textbf{Workload to silicon} \\
\hline
Llama-3.2 1B & Linear, RoPE, MHA, KV\_cache, LayerNorm, Residual, Softmax, Swish, Gate, Sampling & $\sim600$ & $\sim$4 weeks $\times$ person \\
SpinQuant & Dyn. Asym. INT4 / Sta. Sym. INT8 Quant / Dequant, FHT & $\sim200$ & $\sim$1 weeks $\times$ person \\
HMT Plug-in & Linear, MHA, KV\_cache & $\sim200$ & $\sim$1 weeks $\times$ person \\
\hline
\end{tabular}
\label{tab:module_summary}
\vspace{-0.2cm}
\end{table}

\begin{table}[t]
  \centering
  \caption{Ablation Study on Llama-3.2-1B (lower PPL is better). }
  \label{tab:ablation}
  \footnotesize
  \setlength{\tabcolsep}{3pt}  
  \vspace{-0.2cm}
  \renewcommand{\arraystretch}{0.9}
  \resizebox{0.8\linewidth}{!}{
  \begin{tabular}{lcccccccccc}
    \toprule
    \textbf{Quant Config.} & \textbf{W} & \textbf{A} & \textbf{Attn} & \textbf{Vocab} &
    \textbf{Wiki-2 PPL} \\
    \midrule
    No\_Quant           & BF16 & BF16 & BF16        & BF16 & 8.94   \\
    Q0 (SpinQuant)      & INT4 & INT4 & BF16-INT4   & BF16 & 13.30  \\
    Q1                  & INT4 & INT4 & Dyn. INT8    & BF16 & 12.07  \\
    Q2                  & INT4 & INT4 & Sta. INT8    & BF16 & 12.28   \\
    Q3 (Final)          & INT4 & INT4 & Sta. INT8    & INT4 & 12.68 \\
    \bottomrule
  \end{tabular}}
\begin{minipage}{1.0\linewidth}
\scriptsize
\vspace{3pt}
\hspace{15pt}W/A/Attn/Vocab: weight/activation/attention/lm\_head precision.
\end{minipage}
\vspace{-0.8cm}
\end{table}


\begin{table*}[t]
  \centering
  \caption{Parameter configuration of the Llama~3.2-1B model, accelerator, and HMT plug-in, along with the on-board frequency, hardware resource utilization, and latency on the U280 FPGA and synthesis results (estimation marked with~*) on the V80 FPGA.}
  \label{tab:freq-util}
  \vspace{-0.2cm}
  \renewcommand{\arraystretch}{1.1}
  \setlength{\tabcolsep}{3.8pt}
  \resizebox{\textwidth}{!}{
  \begin{tabular}{l l c c c c c c c c}
    \hline
    \textbf{Model/Device} & \textbf{Design} & \textbf{Parameters} & \textbf{Freq.} & \textbf{CLB} & \textbf{DSP} & \textbf{LUT} & \textbf{FF} & \textbf{BRAM / URAM} & \textbf{Latency} \\
    \hline
    Llama 3.2-1B & -- & $L{=}16,\, d{=}2048,\, d_{kv}{=}512,\, d_{ffn}{=}8192,\, d_{lm\_head}{=}128256$ & -- & -- & -- & -- & -- & -- / -- & --\\
    \hline
    \multirow{3}{*}{U280} 
      & Prefill Arch. & $\textit{TP}{=}8,\, \textit{WP}_{\textit{kqvo}}{=}24,\, \textit{WP}_{\textit{mha}}{=}16,\, \textit{WP}_{\textit{ffn}}{=}96$ & 304\,MHz & 66\% & 29\% & 39\% & 24\% & 35\% / 11\% & 1.65\,s/1k tokens\\
      & Decode Arch.  & $\textit{BP}{=}16,\, \textit{WP}_{\textit{int4}}{=}1024,\, \textit{WP}_{\textit{mha}}{=}256$ & 292\,MHz & 76\% & 18\% & 44\% & 28\% & 41\% / 15\% & 6.94\,s/1k tokens\\
      & HMT Plug-in & $N{=}64,\, \textit{BP}{=}4,\, \textit{WP}_{\textit{mem\_attn}}{=}4$ & 290\,MHz & 7.5\% & 1.5\% & 5.3\% & 1.9\% & 4.3\% / 3.8\% & 8.44\,ms/segment\\
    \hline
    \multirow{2}{*}{V80}  
      & Prefill Arch. & $\textit{TP}{=}16,\, \textit{WP}_{\textit{kqvo}}{=}32,\, \textit{WP}_{\textit{mha}}{=}32,\, \textit{WP}_{\textit{ffn}}{=}128$ & 300\,MHz* & 58\% & 26\% & 37\% & 20\% & 22\% / 9\%  & 0.61\,s/1k tokens*\\
      & Decode Arch.  & $\textit{BP}{=}64,\, \textit{WP}_{\textit{int4}}{=}4096,\, \textit{WP}_{\textit{mha}}{=}1024$ & 300\,MHz* & 75\% & 25\% & 42\% & 22\% & 36\% / 20\% & 1.68\,s/1k tokens*\\
      & HMT Plug-in & $N{=}64,\, \textit{BP}{=}4,\, \textit{WP}_{\textit{mem\_attn}}{=}8$ & 300\,MHz* & 3.8\% & 0.7\% & 3.3\% & 0.9\% & 2.4\% / 1.9\% & 6.50\,ms/segment*\\
    \hline
  \end{tabular}}
\vspace{-0.5cm}
\end{table*}

% \subsection{HMT Plug-in Integration}
\section{Case Study 2: HMT Integration} 

% While quantization effectively mitigates computation and memory overheads, long-context prompting still incurs quadratic growth in both compute and memory, resulting in substantial latency and energy consumption in real-world applications. This limitation exposes a major gap in existing LLM accelerator designs. To address this challenge, we integrate the \textit{Hierarchical Memory Transformer (HMT)} mechanism, which recurrently compresses historical context into compact memory embeddings and retrieves relevant information through a memory-attention pathway. The HMT plug-in reduces both computational and memory complexity from quadratic to linear with respect to sequence length.

% Benefiting from the composable and modular design of FlexLLM, we can implement and integrate the HMT plug-in with minimal effort—reusing existing linear and attention modules to form its memory encoder and retrieval pipeline. This seamless integration significantly enhances the accelerator’s capability for efficient and scalable long-context processing while keeping hardware overheads minimal.

% Figure~\ref{fig:Implementation} (c) illustrates the collaborative workflow between the HMT plug-in and the LLM accelerator. Before processing, a long prompt is divided into multiple short segments. In the first stage, the HMT segment processor loads the first half of the current segment ($Seg_n^T$) and concatenates it with a special topic token embedding $T_n$ for segment-level summarization. The combined segment is then passed to the LLM accelerator, which produces an output containing a topic summary vector $S_n$. This $S_n$ is fed back to the HMT module, where it performs cross-attention with the most recent $N$ long-term memory embeddings $\{Mem_{n-N}, \dots, Mem_{n-1}\}$ stored in the memory queue to generate a new prompt embedding $P_n$ that encodes relevant historical information.

% In the next stage, the segment processor augments the current input by concatenating the full segment with $P_n$ and a portion of the previous segment ($Seg_{n-1}^P$) as short-term memory. This enriched segment is sent back to the LLM accelerator to generate a new long-term memory embedding $Mem_n$ for subsequent iterations. 

% By integrating the HMT plug-in, the LLM accelerator reduces its computational and memory complexity from quadratic to linear with respect to the prompt length. This enables efficient long-context processing with minimal hardware overhead and negligible additional latency, significantly extending the applicability of LLM accelerators to real-world, long-form generation tasks.

While quantization reduces compute and memory traffic, long-context prompting still incurs quadratic attention cost in both computation and KV access, leading to high latency and energy. To address this, we integrate a Hierarchical Memory Transformer (HMT) plug-in into our accelerator system. HMT recurrently compresses context segments into compact memory embeddings and retrieves relevant information via a memory-attention pathway, reducing long-context processing complexity from quadratic to linear in sequence length. Moreover, this is an important study to showcase FlexLLM's flexibility and efficiency for implementing new LLM mechanisms.

FlexLLM’s composable design makes this integration lightweight: The HMT plug-in is built largely by reusing existing linear and attention modules to implement its segment encoder, memory generation, and history retrieval pipeline, requiring minimal additional hardware and design effort.

Fig.~\ref{fig:Implementation}(c) shows the end-to-end workflow. A long prompt is split into into multiple short segments. For each segment, the segment processor first forms a summary prompt by concatenating the first half of current segment ($Seg_n^T$) with a topic token embedding $T_n$, and sends it to the LLM accelerator to produce a topic summary vector $S_n$. The HMT plug-in then performs cross-attention between $S_n$ and the most recent $N$ memory embeddings $\{Mem_{n-N},\dots,Mem_{n-1}\}$ in a memory queue to generate a retrieved prompt embedding $P_n$.

Next, the HMT plug-in constructs an augmented prompt by concatenating the full segment, $P_n$, and a short-term context slice from the previous segment ($Seg_{n-1}^P$). This augmented input is processed by the LLM accelerator to produce the new long-term memory embedding $Mem_n$, which is appended to the memory queue for future retrieval.

The experimental results in \ref{results for HMT} show that the HMT plug-in built on FlexLLM enables scalable context window with minimal hardware overhead and negligible added latency over the backbone model implementation, substantially extending the applicability of LLM accelerators to long-context workloads.



\section{Evaluation}

\subsection{Experiment Setup}

\textbf{Hardware Platforms:}  
We select AMD Alveo U280~\cite{amd_u280_ds} at TSMC 16nm and Versal V80~\cite{amd_versal_v80_ds} at 7nm as our target accelerator platforms, and NVIDIA A100 at 7nm as GPU representative.
Table~\ref{tab:hardware_comparison} summarizes key hardware specifications.



% \textbf{Accelerator Implementation.}  
% We implement accelerators with stage-customized architectures for the Llama-3.2 1B model quantized with SpinQuant on both the U280 and V80 FPGAs using FlexLLM. Figure~\ref{fig:Layout} exhibits the design layout on U280 FPGA. In addition, the HMT plug-in is integrated with the accelerator design on the same U280 FPGA to enable efficient long-context processing. Table~\ref{tab:freq-util} summarizes the model configuration, architectural parameters, design frequency, and hardware resource utilization across both platforms. On the U280, the design is deployed on-board and supports rapid architecture reconfiguration within milliseconds, with performance metrics directly measured from real hardware execution. For the V80, we perform RTL synthesis and functional verification using Vitis~2024.2 emulation, and project its expected performance through architectural scaling from the U280 implementation. We choose Allo with W4A8KV8 SmoothQuant as SOTA accelerator enemy.

\textbf{Accelerator Implementations:}  
 Using FlexLLM, we implement stage-customized accelerators for the Llama-3.2 1B model quantized with SpinQuant on both U280 and V80. Fig.~\ref{fig:Layout} shows the U280 layout. We also integrate the HMT plug-in for long-context processing. Table~\ref{tab:freq-util} summarizes model parameters, architectural configuration, achieved frequency, and resource utilization on both platforms. On U280, the design runs on-board with rapid reconfiguration ($\sim$0.3\,s) and performance measured directly on hardware. For V80 implementation, we perform RTL synthesis and P\&R using Vitis/Vivado~2024.2 and AVED toolchain, and project on-board performance by scaling from the U280 implementation. We use Allo~\cite{chen2024allo} with W4A8KV8 SmoothQuant as the SOTA accelerator baseline.

\textbf{GPU Baselines:}
We use the HuggingFace BF16 Llama-3.2 1B model as the GPU baseline. Since vLLM~\cite{kwon2023vllm} does not support SpinQuant, we instead quantize the model to INT4 using GPTQ~\cite{frantar2022gptq} with Marlin~\cite{frantar2025marlin} and evaluate it with vLLM on an A100 GPU. This setup provides a fair GPU reference for our FPGA implementations.

\textbf{Evaluation Metrics:}  
We evaluate FPGA-based accelerators and GPU baselines using three key metrics: (1) end-to-end latency, (2) decode throughput, and (3) energy efficiency. End-to-end latency measures the total inference time (prefill + decode), while decode throughput represents the number of tokens generated per second during the decoding. We record the average power from on-board sampling for the A100 and U280, and synthesis report for the V80,  and compute the tokens-per-joule (token/J) metric to evaluate energy efficiency.

\subsection{Experiment Results}

\begin{figure}[t]
    \centering
    \includegraphics[width=\linewidth]{Figures/U280_Layout.pdf}
    \vspace{-0.8cm}
    \caption{The implementation layout of (a) prefill and (b) decode architectures for quantized Llama-3.2 1B on U280.}
    \label{fig:Layout}
    \vspace{-0.6cm}
\end{figure}

\begin{figure*}[t]
    \centering
    \includegraphics[width=1.0\linewidth]{Figures/Pref_Norm_new.pdf}
    \vspace{-0.8cm}
    \caption{Performance and energy efficiency comparison of FlexLLM-based accelerators and A100 GPU under standard inference settings.}
    \label{fig:Perf_Comparison}
    \vspace{-0.2cm}
\end{figure*}


\begin{figure*}[t]
    \centering
    \includegraphics[width=1.0\linewidth]{Figures/Perf_HMT_new.pdf}
    \vspace{-0.9cm}
    \caption{Performance and energy efficiency comparison in long-context scenarios with and without the HMT plug-in.}
    \label{fig:Perf_HMT}
    \vspace{-0.5cm}
\end{figure*}


\subsubsection{Case Study 1: Typical Sequence Lengths}  
% We evaluate the proposed accelerators and the GPU baseline across various prefill and decode sequence lengths, as illustrated in Figure~\ref{fig:Perf_Comparison}. As expected, the GPU demonstrates a clear advantage in scenarios with long prefill and short decode lengths (e.g., [1024, 256]), where its massive parallel compute capability yields ultra-fast prompt ingestion. However, as the decode length increases, the performance bottleneck of the GPU shifts to memory bandwidth, whereas our FPGA-based accelerators maintain high efficiency through stage-customized module designs and flexible parallelism.  

We evaluate the proposed accelerators and the GPU baseline across a range of prefill and decode sequence lengths (Fig.~\ref{fig:Perf_Comparison}). As expected, the GPU demonstrates a clear advantage in scenarios with long prefill and short decode lengths (e.g., [1024, 256]), where its massive parallel compute capability yields ultra-fast prompt ingestion. However, as decode length grows, the performance bottleneck of the GPU shifts to memory bandwidth, while our FPGA accelerators sustain higher efficiency via stage-customized designs and flexible parallelism.

% On average, the FlexLLM-based accelerator on the U280 achieves \textbf{1.29$\times$} end-to-end speedup and \textbf{1.64$\times$} higher decode throughput than the A100 BF16 baseline, surpassing Allo with \textbf{1.46$\times$} and \textbf{1.35$\times$}, respectively. The accelerator on V80 further delivers up to \textbf{4.71$\times$} end-to-end speedup and \textbf{6.55$\times$} higher decode throughput. For longer decoding sequences (e.g., [512, 2048] and [1024, 2048]), the design on V80 outperforms the A100 running the INT4 GPTQ-Marlin model under vLLM Although A100 has higher memory bandwidth, our accelerator can still outperform in memory-bound decode stage due to higher bandwidth utilization, where A100 with vLLM only achieve 13.06\% average memory utilization, our on V80 achieve average 38.61\% and peak up to 93.65\%. This demonstrating the superior performance and efficiency of our FlexLLM-enabled stage-customized architecture. 

% In terms of energy efficiency, the FlexLLM-based accelerators on both U280 and V80 show clear advantages over the A100 BF16 baseline. On average, the U280 achieves \textbf{3.14$\times$} and the V80 \textbf{4.13$\times$} higher tokens-per-joule. When the decode length increases to 1024 or 2048 tokens, even the 16\,nm U280 surpasses the 7\,nm A100 running GPTQ-Marlin with vLLM, demonstrating superior resource utilization and power efficiency of our design. Compared to Allo, our accelerator also achieves \textbf{1.10$\times$} higher energy efficiency. These results confirm that FlexLLM enables domain-specific accelerators that outperform general-purpose GPUs in both performance and energy efficiency, even across process nodes.

On average, the FlexLLM-based accelerator on U280 achieves \textbf{1.29$\times$} end-to-end speedup and \textbf{1.64$\times$} higher decode throughput over the A100 BF16 baseline, and also surpasses Allo with \textbf{1.46$\times$} and \textbf{1.35$\times$}, respectively. The V80 design further delivers up to \textbf{4.71$\times$} end-to-end speedup and \textbf{6.55$\times$} higher decode throughput. For longer decode workloads (e.g., [512, 2048] and [1024, 2048]), V80 even outperforms A100 running INT4 GPTQ-Marlin under vLLM. Although A100 has higher peak bandwidth, decode performance is bounded by effective bandwidth utilization: A100+vLLM achieves only 13.06\% average bandwidth utilization, whereas our V80 design reaches \textbf{38.61\%} on average and up to \textbf{74.93\%} peak. These results highlight the benefit of FlexLLM-enabled stage-customized architectures for memory-bound decode.

In energy efficiency, FlexLLM-based accelerators outperform the A100 BF16 baseline, achieving \textbf{3.14$\times$} (U280) and \textbf{4.13$\times$} (V80) higher tokens-per-joule on average. With long decode lengths (1024/2048), even the 16\,nm U280 surpasses the 7\,nm A100 running INT4 GPTQ-Marlin with vLLM, demonstrating strong utilization and power efficiency. Compared to Allo, our accelerator further achieves \textbf{1.10$\times$} higher energy efficiency. These results confirm that FlexLLM enables domain-specific accelerators that outperform general-purpose GPUs in both performance and energy efficiency, even across process nodes.




\subsubsection{Case Study 2: Long Context with HMT Plug-in}
\label{results for HMT}

% In long-context inference, accelerators without the HMT plug-in face a high risk of running out of memory due to the limited HBM capacity on FPGAs. Even when sufficient off-chip memory is available for KV caching, the theoretical prefill latency on the U280 can exceed one hour when the input sequence reaches 64K tokens—clearly impractical for real-world deployment. These results underscore the necessity of sequence compression mechanisms specifically designed for LLM accelerators, such as the HMT plug-in.

% According to Table~\ref{tab:freq-util}, the HMT plug-in implemented with FlexLLM consumes less than 7.5\% of total hardware resources and introduces only 8.44\,ms additional latency per segment—negligible (\textbf{\textless0.6\%}) compared with the overall accelerator latency. Moreover, thanks to FlexLLM’s templated modular design, users can flexibly tune the plug-in’s resource footprint and processing delay to meet different backbone accelerator integration requirements.

% Figure~\ref{fig:Perf_HMT} evaluates the performance and energy efficiency of accelerators enhanced with the HMT plug-in. Without HMT, the A100 GPU still outperforms the FPGA-based accelerators in the prefill stage due to its massive compute throughput. However, when compared to the theoretical latency bound without HMT (Equation~\ref{eq:prefill_latency}), integrating HMT achieves up to a \textbf{23.23$\times$} reduction in prefill latency—bringing the time-to-first-token (TTFT) to a practically acceptable range for long-context tasks. 

% For end-to-end performance, HMT-enhanced accelerators regain a clear advantage over the A100, driven by their superior decode throughput. This improvement stems from two factors: (1) the accelerator’s inherently efficient compute and memory architecture, and (2) the HMT plug-in’s sequence compression, which transforms decode-stage complexity from quadratic to linear scaling with sequence length. As a result, the HMT-augmented design achieves up to \textbf{1.36$\times$} higher decode throughput and \textbf{1.10$\times$} lower end-to-end latency on the U280 compared to the A100 baseline, and up to \textbf{4.86$\times$} and \textbf{3.70$\times$} improvements, respectively, on the V80.


% In terms of energy efficiency, the HMT plug-in helps maintain a stable tokens-per-joule rate as the sequence length grows, while the A100’s efficiency degrades sharply in long-context scenarios. Consequently, the HMT-enhanced accelerators on U280 and V80 exhibit substantial advantages, achieving \textbf{5.21$\times$} and \textbf{6.27$\times$} higher energy efficiency than the A100 baseline, and \textbf{1.65$\times$} and \textbf{1.99$\times$} improvements over the A100 optimized with GPTQ-Marlin and vLLM. These results highlight not only the critical role of the HMT plug-in in enabling scalable long-context inference, but also the flexibility and extensibility of our FlexLLM framework for domain-specific LLM accelerator design.

In long-context inference, accelerators without HMT are prone to memory exhaustion due to limited FPGA HBM capacity. Even if off-chip memory suffices for KV caching, the theoretical prefill latency on U280 can exceed one hour at 64K tokens, making deployment impractical. This motivates sequence compression tailored for LLM accelerators, such as HMT.

As shown in Table~\ref{tab:freq-util}, the HMT plug-in on U280 consumes under \textbf{7.5\%} of total resources and adds only 8.44\,ms per segment (\textbf{\textless0.6\%} of end-to-end latency). FlexLLM’s templated modularity further enables tuning the plug-in’s resource--latency trade-off to match different LLM accelerators.

Figure~\ref{fig:Perf_HMT} reports performance and energy efficiency with HMT. The A100 without HMT remains faster than FPGA-based accelerators in prefill due to higher peak throughput; however, compared to the theoretical latency bound without HMT, HMT reduces prefill latency by up to \textbf{23.23$\times$}, bringing the time-to-first-token into a practical range for long prompts. For end-to-end performance, HMT-enhanced accelerators regain advantage over the A100, driven by higher decode throughput and HMT’s compression that converts the decode scaling from quadratic to linear in sequence length. Compared with the A100 BF16 baseline, the U280 achieves up to \textbf{1.36$\times$} higher decode throughput and \textbf{1.10$\times$} lower end-to-end latency, while the V80 reaches up to \textbf{4.86$\times$} and \textbf{3.70$\times$}, respectively.

HMT also stabilizes energy efficiency as context grows: tokens/J drops sharply on A100 for long prompts, whereas the HMT-enhanced U280 and V80 sustain high efficiency, delivering up to \textbf{5.21$\times$} and \textbf{6.27$\times$} higher tokens/J than the A100 BF16 baseline, and \textbf{1.65$\times$} and \textbf{1.99$\times$} over GPTQ-Marlin+vLLM. These results highlight not only the critical role of the HMT plug-in in enabling scalable long-context inference, but also the flexibility and extensibility of our FlexLLM framework for efficient LLM accelerator design.







\section{Conclusion and Outlook}
We present \textbf{FlexLLM}, a composable HLS framework for rapidly building and optimizing domain-specific LLM accelerators. FlexLLM enables stage-customized hybrid architectures for prefill and decode, and provides a comprehensive quantization stack for accurate low-bit inference, bridging model-level innovation and hardware-level implementation. Experimental results on the Llama-3.2 1B model demonstrate that FlexLLM enables high-performance and energy-efficient LLM inference designs, achieving performance comparable to or surpassing the NVIDIA A100 GPU running with vLLM. These results underscore FlexLLM’s capability to translate algorithmic innovations into efficient hardware realizations with minimal manual effort. Looking ahead, FlexLLM provides a foundation for domain-specific LLM acceleration on future platforms and for extending composable HLS methodologies to next-generation Transformer models.




% conference papers do not normally have an appendix


% use section* for acknowledgment
% \section*{Acknowledgment}


% The authors would like to thank...





% trigger a \newpage just before the given reference
% number - used to balance the columns on the last page
% adjust value as needed - may need to be readjusted if
% the document is modified later
%\IEEEtriggeratref{8}
% The "triggered" command can be changed if desired:
%\IEEEtriggercmd{\enlargethispage{-5in}}

% references section

% can use a bibliography generated by BibTeX as a .bbl file
% BibTeX documentation can be easily obtained at:
% http://mirror.ctan.org/biblio/bibtex/contrib/doc/
% The IEEEtran BibTeX style support page is at:
% http://www.michaelshell.org/tex/ieeetran/bibtex/
\bibliographystyle{IEEEtran}

% argument is your BibTeX string definitions and bibliography database(s)
%\bibliography{IEEEabrv,../bib/paper}
\bibliography{example_paper}
%
% <OR> manually copy in the resultant .bbl file
% set second argument of \begin to the number of references
% (used to reserve space for the reference number labels box)
% \begin{thebibliography}{1}

% \bibitem{IEEEhowto:kopka}
% H.~Kopka and P.~W. Daly, \emph{A Guide to \LaTeX}, 3rd~ed.\hskip 1em plus
%   0.5em minus 0.4em\relax Harlow, England: Addison-Wesley, 1999.

% \end{thebibliography}




% that's all folks
\end{document}


